\documentclass[11pt,reqno]{amsart}

\setlength{\topmargin}{0in}
\setlength{\textheight}{8.6in}
\setlength{\oddsidemargin}{0in}
\setlength{\evensidemargin}{0in}
\setlength{\textwidth}{6.5in}
\usepackage[usenames,dvipsnames]{xcolor}
\usepackage[normalem]{ulem}
\usepackage{amsmath,amssymb,amsthm,graphics,amscd,mathrsfs}
\usepackage{longtable}
\usepackage{cite}
\usepackage[all,cmtip]{xy}
\usepackage{graphicx,wrapfig}
\usepackage{epsf}

\usepackage{hyperref}
\hypersetup{colorlinks=true,citecolor=Purple}


\usepackage{lscape}
\setlength{\unitlength}{1mm}
\usepackage{tikz}
\usetikzlibrary{matrix,arrows.meta}

\usepackage{verbatim}


\renewcommand{\thesection}{\Roman{section}}

\usepackage{etoolbox}

% \DeclareRobustCommand\skipifequals[2]{%
%   \ifnum\pdfstrcmp{#1}{#2}=0
%   \else
%     #2%
%   \fi
% }
% \newcommand\skipifcurrent[1]{\protecting{\skipifequals{#1}}{#1}}
% \renewcommand\thesubsection{\skipifcurrent{\thesection.}\arabic{subsection}}
% \renewcommand\thesubsubsection{\skipifcurrent{\thesection.}\arabic{subsection}.\arabic{subsubsection}}


\usepackage{titlesec}
\titleformat{\section}{\Large\bfseries\center\sc}
{Part\ \thesection\,---}{0.3ex}{} %% Text "Lecture" is what you would like the section number to display with.
%\titlespacing{\subsection}{2em}{*1}{*1}
\titleformat{\subsection}{\large\sc\center}{\thesubsection.}{0.5ex}{}
\titleformat{\subsubsection}[runin]{\bf}{\thesubsubsection.}{0.5ex}{}[.]

\newcommand{\asection}[2]{
\setcounter{section}{#1}
\addtocounter{section}{-1}
\subsubsection{#2}
}

\setcounter{tocdepth}{4}
\setcounter{secnumdepth}{3}



\let\oldthebibliography=\thebibliography
  \let\endoldthebibliography=\endthebibliography
  \renewenvironment{thebibliography}[1]{
    \begin{oldthebibliography}{#1}
      \setlength{\parskip}{0ex}
      \setlength{\itemsep}{0ex}
  }
  {
    \end{oldthebibliography}
  }
  
\setlongtables


\renewcommand{\labelenumi}{(\roman{enumi})}

\usepackage[capitalise]{cleveref}
\crefformat{equation}{#2(#1)#3}

\crefformat{enumi}{#2(\romannumeral #1\relax)#3}
\crefformat{enumii}{#2(\romannumeral\theenumi\relax)(\theenumii)#3}


\newtheorem{lemma}{Lemma}[subsection]
\newtheorem{theorem}[lemma]{Theorem}
\newtheorem*{claim}{Claim}
\newtheorem{cor}[lemma]{Corollary}
\newtheorem{conjecture}[lemma]{Conjecture}
\newtheorem{prop}[lemma]{Proposition}
\newtheorem{question}[lemma]{Question}
\theoremstyle{definition}
\newtheorem{example}[lemma]{Example}
\newtheorem{examples}[lemma]{Examples}
\newtheorem{algorithm}[lemma]{Algorithm}
\newtheorem*{algorithm*}{Algorithm}
\theoremstyle{remark}
\newtheorem{remark}[lemma]{Remark}
\newtheorem{remarks}[lemma]{Remarks}
\newtheorem*{remarks*}{Remarks}
\newtheorem{obs}[lemma]{Observation}
\theoremstyle{definition}
\newtheorem{defn}[lemma]{Definition}



\theoremstyle{plain}
\newtheorem{theorem1}{Theorem}%[blah]
\crefformat{theorem1}{Theorem #2#1#3}



\makeatletter
\newcommand{\pushright}[1]{\ifmeasuring@#1\else\omit\hfill$\displaystyle#1$\fi\ignorespaces}
\newcommand{\pushleft}[1]{\ifmeasuring@#1\else\omit$\displaystyle#1$\hfill\fi\ignorespaces}
\makeatother

   
  \def\hal{\unskip\nobreak\hfil\penalty50\hskip10pt\hbox{}\nobreak
  \hfill\vrule height 5pt width 6pt depth 1pt\par\vskip 2mm}


\DeclareMathOperator{\Lie}{Lie}
\DeclareMathOperator{\im}{im}
\DeclareMathOperator{\Rad}{Rad}
\DeclareMathOperator{\cl}{cl}
\DeclareMathOperator{\gr}{gr}
\DeclareMathOperator{\rk}{rk}
\DeclareMathOperator{\id}{id}
\DeclareMathOperator{\ad}{ad}
\DeclareMathOperator{\Dist}{Dist}
\DeclareMathOperator{\Spec}{Spec}
\DeclareMathOperator{\Der}{Der}
\DeclareMathOperator{\Char}{char}
\DeclareMathOperator{\Res}{Res}
\DeclareMathOperator{\Soc}{Soc}
\DeclareMathOperator{\Ind}{Ind}
\DeclareMathOperator{\Gal}{Gal}
\DeclareMathOperator{\Mat}{Mat}
\DeclareMathOperator{\Ann}{Ann}
\DeclareMathOperator{\Jac}{Jac}
\DeclareMathOperator{\In}{in}
\DeclareMathOperator{\Gr}{Gr}

\newcommand{\len}{\mathit{l}}
\newcommand{\x}{$\times$}
\newcommand{\bb}{\mathbf}
\newcommand{\R}{\mathrm{R}}
\newcommand{\RR}{\mathscr{R}}
\newcommand{\G}{\mathscr{G}}
\newcommand{\GG}{\mathcal{G}}
\newcommand{\TT}{\mathcal{T}}
\renewcommand{\H}{\mathscr{H}}
\newcommand{\hra}{\hookrightarrow}
\newcommand{\sss}{\mathrm{ss}}
\newcommand{\Hom}{\mathrm{Hom}}
\newcommand{\Int}{\mathrm{int}}
\newcommand{\Ext}{\mathrm{Ext}}
\newcommand{\Mor}{\mathrm{Mor}}
\newcommand{\opH}{\mathrm{H}}
\newcommand{\D}{\mathscr{D}}
\newcommand{\SO}{\mathrm{SO}}
\newcommand{\Sp}{\mathrm{Sp}}
\newcommand{\SL}{\mathrm{SL}}
\newcommand{\SU}{\mathrm{SU}}
\newcommand{\GL}{\mathrm{GL}}
\newcommand{\kp}{k_{\mathbf{p}}}
\newcommand{\PGL}{\mathrm{PGL}}
\newcommand{\OO}{\mathcal{O}}
\newcommand{\Y}{\mathbf{Y}}
\newcommand{\X}{\mathbf{X}}
\newcommand{\diag}{\mathrm{diag}}
\newcommand{\End}{\mathrm{End}}
\newcommand{\tr}{\mathrm{tr}}
\newcommand{\Stab}{\mathrm{Stab}}
\newcommand{\red}{\mathrm{red}}
\newcommand{\Aut}{\mathrm{Aut}}
\renewcommand{\H}{\mathcal{H}}
\renewcommand{\u}{\mathfrak{u}}
\newcommand{\Ad}{\mathrm{Ad}}
\newcommand{\N}{\mathcal{N}}
\newcommand{\NN}{\mathbb{N}}
\newcommand{\Z}{\mathbb{Z}}
\newcommand{\la}{\langle}\newcommand{\ra}{\rangle}
\newcommand{\gl}{\mathfrak{gl}}
\newcommand{\g}{\mathfrak{g}}
\newcommand{\F}{\mathbb{F}}
\newcommand{\m}{\mathfrak{m}}
\renewcommand{\b}{\mathfrak{b}}
\newcommand{\p}{\mathfrak{p}}
\newcommand{\q}{\mathfrak{q}}
\renewcommand{\l}{\mathfrak{l}}
\newcommand{\del}{\partial}
\newcommand{\h}{\mathfrak{h}}
\renewcommand{\t}{\mathfrak{t}}
\renewcommand{\k}{\mathfrak{k}}
\newcommand{\Gm}{\mathbb{G}_m}
\renewcommand{\c}{\mathfrak{c}}
\renewcommand{\r}{\mathfrak{r}}
\newcommand{\n}{\mathfrak{n}}
\newcommand{\s}{\mathfrak{s}}
\newcommand{\Q}{\mathbb{Q}}
\newcommand{\z}{\mathfrak{z}}
\newcommand{\pso}{\mathfrak{pso}}
\newcommand{\so}{\mathfrak{so}}
\renewcommand{\sl}{\mathfrak{sl}}
\newcommand{\psl}{\mathfrak{psl}}
\renewcommand{\sp}{\mathfrak{sp}}
\newcommand{\Ga}{\mathbb{G}_a}
\newcommand{\barB}{\overline{B}}
\newcommand{\barm}{\overline{\mathfrak{m}}}
\newcommand{\CC}{\mathbb{C}}
\newcommand{\KK}{\mathcal{K}}
\newcommand{\Cs}{\mathsf{C}}
\newcommand{\ks}{\mathsf{k}}
\newcommand{\Ts}{\mathsf{T}}
\newcommand{\CCC}{\mathcal{C}}
\newcommand{\pred}{\mathrm{pred}}
\newcommand{\sep}{\mathrm{sep}}

\newcommand{\B}{\mathscr{B}}
\newcommand{\C}{\mathscr{C}}
\renewcommand{\D}{\mathscr{D}}
\newcommand{\U}{\mathscr{U}}
\newcommand{\V}{\mathscr{V}}

\newcommand{\bark}{{\overline{k}}}
\newcommand{\barG}{{\overline{G}}}
\newcommand{\Ol}[1]{{\color{OliveGreen} #1}}


\setlength{\parskip}{1mm}

\title{Geometric rigidity of simple modules for algebraic groups}
\subjclass[2010]{20G15}

\author[Bate]{Michael Bate}
\address
{Department of Mathematics,
University of York,
York YO10 5DD,
United Kingdom}
\email{michael.bate@york.ac.uk}
\author[Stewart]{David I. Stewart*}
    \address{Department of Mathematics, The University of Manchester, Manchester, UK}
    \email{david.i.stewart@manchester.ac.uk}

\begin{document}

\begin{abstract}
Let $k$ be a field, let $G$ be a smooth affine $k$-group and $V$ a finite-dimensional $G$-module. 
We say $V$ is \emph{rigid} if the socle series and radical series coincide for the action of $G$ on each indecomposable summand of $V$; say $V$ is \emph{geometrically rigid} (resp.~\emph{absolutely rigid}) if $V$ is rigid after base change of $G$ and $V$ to $\bar k$ (resp.~any field extension of $k$).
We show that all simple $G$-modules are geometrically rigid, though not in general absolutely rigid.
More precisley, we show that if $V$ is a simple $G$-module, then there is a finite purely inseparable extension $k_V/k$ naturally attached to $V$ such that
$V_{k_V}$ is absolutely rigid as a $G_{k_V}$-module. 
The proof for connected $G$ turns on an investigation of algebras of the form $K\otimes_k E$ where $K$ and $E$ are field extensions of $k$; 
we give an example of such an algebra which is not rigid as a module over itself. 
We establish the existence of the purely inseparable field extension $k_V/k$ through an analogous version for artinian algebras.

In the second half of the paper we apply recent results on the structure and representation theory of pseudo-reductive groups to gives a concrete description of $k_V$. Namely, we combine  the main structure theorem of the Conrad--Prasad classification of pseudo-reductive $G$ together with our previous high weight theory. For $V$ a simple $G$-module, we calculate the minimal field of definition of the geometric Jacobson radical of $\End_G(V)$ in terms of the high weight of $G$ and the Conrad--Prasad classification data; this gives a concrete construction of the field $k_V$ as a subextension of the minimal field of definition of the geometric unipotent radical of $G$. 

 We also observe that the Conrad--Prasad classification can be used to hone the dimension formula for $G$ we had previously established; we also use it to give a description of $\End_G(V)$ which includes a dimension formula.
\end{abstract}

\maketitle

\section*{Introduction}

Let $k$ be a field and $G$ a smooth affine $k$-group.
The recent classification by highest weight of the (rational) simple $G$-modules for connected $G$ in \cite{BS22} has opened the possibility of answering general
questions about the representation theory of algebraic groups, which hitherto might have seemed inaccessible; this paper is presented in that spirit. 
Given a simple $G$-module $V$ we provide rather detailed information about the behaviour of $V$ under field extensions---we describe the structure of $V_{E}$ as a $G_{E}$-module for suitable field extensions $E/k$. 
Principally, we prove that whilst $V$ is far from being absolutely simple in general, or even absolutely semisimple, it is at least \emph{geometrically rigid}. For any finite-dimensional $G$-module $V$, we say 
that $V$ is \emph{rigid} if the socle series and radical series coincide for the action of $G$ on each of its indecomposable summands; we say $V$ is \emph{geometrically rigid} (resp.~\emph{absolutely rigid}) if $V$ is rigid after base change to $\bark$ (resp.~any field extension of $k$). Explicit definitions of the above are to be found in \cref{sec:soc}.


Our main result is
\begin{theorem1}\label{thm:mainone}
Let $G$ be a smooth affine $k$-group and $V$ a simple $G$-module.
Then there exists a finite purely inseparable extension $k_V/k$ naturally attached to $V$ such that after base-change to $k_V$, the $G_{k_V}$-module $V_{k_V}$ is absolutely rigid.
In particular, $V$ is geometrically rigid.
\end{theorem1}

\begin{remarks*}
  \begin{enumerate}\item The notion of rigidity has mostly been investigated in the context of indecomposable modules and there is a choice as to how best to generalise it. The obvious alternative would be to say $V$ is \emph{rigid\,$'$} if its socle and radical series coincide; however, this would rule out the direct sums of rigid$'$ modules being rigid$'$ unless they had the same Loewy length. Since many algebraic properties are defined locally---such being Gorenstein---our definition seems more sensible. In particular, the category of rigid modules is additive.
    
\item One of the earliest results about rigidity is due to Jennings: over a field, the group algebra of a finite $p$-group is rigid---\!\cite[p93]{Ben98}. 

\item Obviously any simple module $V$ is rigid, but even when $G$ is connected and smooth, $V$ is not in general absolutely semisimple or absolutely indecomposable, so a statement about how rigidity behaves under base change is not immediate---see \cref{rem:notabs} below. 
It turns out that one can find examples $(G,V,L/E/k)$, where $G$ is a smooth connected $k$-group, $V$ is a simple $G$-module, and $L/E/k$ is a tower of finite extensions, such that $V_E$ is \emph{not} rigid but $V_L$ is absolutely rigid. 
This boils down to exhibiting a tensor products of finite purely inseparable field extensions $K$ and $E$ of $k$ such that  the algebra $A:=K\otimes_k E$ is not rigid as a module over itself---see \cref{ex:finally}. 

\item A natural source of indecomposable non-simple modules for a split reductive group $G$ over a field $k$ are Weyl modules $V(\lambda)$ and tilting modules $T(\lambda)$, where we refer the reader to \cite[II.2.13,\ II.8.3,\ E.3]{Jan03} for definitions. 
It is natural to ask when these modules are rigid. This question is given a thorough treatment in \cite{And11}, and one finds that for characteristic sufficiently large, both Weyl modules and tilting modules are indeed rigid. Non-rigid examples of both are provided in \emph{op.~cit.} as are more examples, again for $\SL_3$, in \cite{BDM} when $p=3$. See also \cite{Haz17} for a novel approach to this problem.
\end{enumerate}
\end{remarks*}


We divide into two parts. Part \ref{parti} is dedicated to the proof of \cref{thm:mainone} and we exhibit an appropriate field $k_V$ via a rather general argument
involving the behaviour of finite-dimensional simple $k$-algebras under base change.
Given such a $k$-algebra $A$, we show that the Jacobson radical $\Jac(A_\bark)$ of the base change of $A$ to $\bark$ has a descent to a minimal field of definition $k'$ where $k'/k$ is purely inseparable; moreover, $A_{k'}$ is absolutely rigid; see \cref{thm:Drigid}.
This result relates to \cref{thm:mainone} when we let $A=\End_G(V)$ be the endomorphism ring of the simple $G$-module $V$ and define $k_V$ to be the minimal field of definition of $\Jac(A_\bark)$---by a Morita equivalence this is the same as the minimal field of definition of $\Jac(Z(A_\bark))=\Jac(Z(A)\otimes_k\bark)$.
During this part of the paper, we also construct an example of two finite purely inseparable extensions $K$ and $E$ such that the regular module of $K\otimes_k E$ is not rigid and deduce a large class of $G$-modules which are not absolutely rigid.  We then extend our analysis to non-connected $G$, completing the proof of \cref{thm:mainone}.

Part \ref{partii} sharpens the conclusion of \cref{thm:mainone} using the high weight theory of \cite{BS22} for pseudo-reductive groups together with the Conrad--Prasad structure theorem describing their classification, 
\cite[Thm.~9.2.1]{CPClass}. We consider endomorphism rings of simple modules in \cref{sec:endsimple1,sec:endsimple2}, whose deliberations afford a concrete construction of $k_V$.
Essentially this reduces to the case where $G$ is a pseudo-split pseudo-reductive group with no non-trivial normal unipotent subgroup, hence locally of minimal type, hence described by the Conrad--Prasad structure theorem. (Recall that $G$ is \emph{pseudo-split} if it has a split maximal torus.) Then by \cite{BS22}, a simple module $V$ is isomorphic to $L_G(\lambda)$ and the field $k_V$ we use coincides with the endomorphism algebra $\End_G(V)$, which also identifies with the high-weight space $L_G(\lambda)_\lambda$; it can be precisely described as a compositum of purely inseparable field extensions using arithmetic information about $\lambda$ together with the the Conrad--Prasad data defining $G$. Mostly, the root system of $G$ has no bearing on $k_V$, while evidently it does on $L_G(\lambda)$. 
In case $G$ is an arbitrary connected smooth affine algebraic $k$-group, we elucidate the structure of the division algebra $D:=\End_G(V)$. We show that $D$ has a unique \emph{$p$-splitting field}: there is a unique minimal extension $E/k$ such that $D\otimes_k E$ is a product of matrix algebras over purely inseparable extensions of $E$. 
Thus engaged, we interpret our previous dimension formula in terms of the same data, and give a formula for the dimension of $D$: \cref{cor:dimforsimples}.

A useful auxiliary result locates the simple modules for pseudo-split pseudo-reductive groups as submodules of simple modules for Weil restrictions of reductive groups. More specifically, for a pseudo-reductive group $G$, there is a homomorphism $i_G:G\to \R_{k_V/k}(G')$, where $k_V$ is the minimal field of definition for the geometric unipotent radical of $G$; the group $G'$ is the corresponding reductive quotient of $G_{k_V}$; and $\R_{k_V/k}$ denotes the Weil restriction functor. We show that when $G'$ is pseudo-split, the simple modules for $\R_{k_V/k}(G')$ are semisimple and isotypic upon restriction to the image of $G$; see \cref{prop:splitend}.

\section[{Proof of \texorpdfstring{\cref{thm:mainone}}{Theorem 1}}]{Proof of \texorpdfstring{\cref{thm:mainone}}{Theorem 1}}\label{parti}
\subsection{Preliminaries}\label{sec:prelims}


Our main references for the theory of algebraic groups are \cite{CGP15}, \cite{Milne17}, and \cite{Jan03}, with the last also our standard reference for
the representation theory of algebraic groups. 
In most of the paper, $k$ denotes a field, but below we do need to consider the base change of $k$-groups and modules to
more general $k$-algebras so up until \cref{cor:DEZE} we also let $k$ denote a general commutative unital ring. 
For such a $k$, we view an affine $k$-group $G$ as a functor $\underline{k\textrm{-Alg}}\to \underline{\mathrm{Grp}}$
which is represented by a $k$-algebra $k[G]$; in other words $G(?)\cong \Hom_{k\textrm{-Alg}}(k[G],?)$. Since the term \emph{algebraic} has a wide range of uses we offer the following clarifications: when $G^0$ makes sense (such as when $k$ is artinian) and $k[G^0]$ is finitely generated, then we say $G$ is \emph{locally of finite type}; if $k[G]$ is finitely generated then we say $G$ is of \emph{finite type}. More restrictively, if $k[G]$ (resp.~$k[G^0]$) is isomorphic to $k[T_1,\dots,T_n]/I$ for $I$ finitely generated, then we say $G$ is \emph{finitely presented} or \emph{algebraic} (resp.~\emph{locally finitely presented}). (If $k$ is artinian then $k[G]$ being finitely generated and locally finitely presented coincides with being finitely presented.)

\begin{center}\framebox{In what follows $G$ will always denote an affine algebraic $k$-group scheme.}\end{center}

In particular, suppose $k$ is a field, $k_s$ its separable closure, and $\bark$ its algebraic closure.
If $G$ is smooth, then $G$ is geometrically reduced, and it follows from \cite[Cor. 1.17]{Milne17}
that $G_{k_s}(k_s)$ is dense in $G_{k_s}$.

\subsubsection{Modules for algebraic groups}\label{sec:Gmods} 
Let $k$ be a commutative unital ring.
Let $M$ be a $k$-module (possibly not finitely generated). 
Then we may define a group functor $M_a:\underline{k\textrm{-Alg}}\to \underline{\mathrm{Grp}}$ so that 
$M_a(A)=M\otimes_k A$ inherits a group structure from the additive group on $A$. 
Note that, even when $k$ is a field, $M_a$ is only an algebraic group when $M$ is finite-dimensional. 
Recall that a left action of $G$ on a $k$-functor $X$ is a morphism (i.e. a natural transformation) 
$\phi:G\times X\to X$ such that $\phi(A):G(A)\times X(A)\to X(A)$ is a left action of the group $G(A)$ on $X(A)$ for each $k$-algebra $A$. 
In case $G$ acts on $M_a$ such that the action of $G(A)$ on $M(A)$ is $A$-linear for each $k$-algebra $A$, we say $M$ is a \emph{representation for $G$}, or more frequently in this paper, a \emph{$G$-module}. 
These definitions allow us to work with arbitrary $k$-modules, 
although in the case $M\cong k^n$ is a $G$-module 
then it corresponds to a homomorphism $G\to \GL(M)$ of algebraic groups. 
If $k\to E$ is a homomorphism of rings and $M$ is a $G$-module then $M_{E}:=M\otimes E$ acquires an action of the base change $G_{E}$ of $G$ making it into a 
$G_{E}$-module.
For more on these definitions, see \cite[\S I.2]{Jan03}.

Morphisms between $G$-modules are $G$-equivariant $k$-linear maps. 
If $M$ and $N$ are $G$-modules, the full collection of such morphisms is written $\Hom_{G}(M,N)$. If $M=N$, we write instead $\End_G(M)$. If $G$ is flat---which is to say that $k[G]$ is a flat $k$-module---then the category of $G$-modules is abelian; i.e. kernels and cokernels are submodules \cite[I.2.9]{Jan03}. Therefore, we have the following, with the usual proof.
\begin{lemma}[Schur] Let $G$ be flat and let $M$ be a simple $G$-module. Then $\End_G(M)$ is a division ring.\label{lem:schur}\end{lemma}
Recall a $G$-module $M$ is \emph{locally finite} if any element $m\in M$ is contained in a $G$-submodule of $M$ which is finitely generated as a $k$-module. It has been noticed by Wilberd van der Kallen \cite[\S1.7]{vdk21} that the proof in \cite[I.2.13]{Jan03} that $G$-modules are locally finite for arbitrary $k$ and flat $G$ is incomplete. It relies on the assumption that an arbitrary intersection of $G$-submodules is again a submodule; however, there is a counterexample to be found at \cite[Expos\'e VI, Edition 2011, Remarque 11.10.1]{SGA3}. Local finiteness is  used to infer finite generation over $k$ of a simple $G$-module, which means that the proof of \cite[I.10.15]{Jan03} is also incomplete. If $M$ is a finitely presented or projective module over $k$ then it commutes with arbitrary direct products, since free modules do. From this is follows that the functor $N\mapsto M\otimes_k N$ commutes with arbitrary direct products, hence with arbitrary limits, of which arbitrary intersections are one example. Therefore if $k$ is noetherian or $k[G]$ is projective, then for any $v\in N$ we can find a finitely generated $G$-submodule $kGv$ of $N$ as the intersection of all those containing $v$ and local finiteness holds.
%However, if $M$ is a module which is either projective or finitely presented over $k$ then the functor $N\mapsto M\otimes_k N$ commutes with arbitrary intersections. This follows because the natural map \[M\otimes_k\prod N_i\to \prod M\otimes_k N_i\] is injective for such $M$; certainly  which is seen from the definition of the tensor product by using direct products. It follows that if  $k$ is noetherian or $k[G]$ is a projective $k$-module, then $G$-modules \emph{are} all locally finite, so that simple modules are  finitely generated over $k$.

\begin{remark}It would be lovely to know if there is an example of a non-flat $G$ with a simple $M$ and non-invertible non-zero endomorphism. Similarly, it would be lovely to know if local finiteness holds more generally than under the hypotheses above, even if only as far as to establish that simple $G$-modules are finitely generated over $k$.\end{remark}
For a flat $k$-algebra $E$ (e.g., an extension of fields $E/k$), we have \cite[I.2.10(7)]{Jan03}
\begin{equation}\label{eq:tensorhom}
\Hom_G(M,N)\otimes_k E \cong \Hom_{G_{E}}(M_{E},N_{E}). 
\end{equation}




\subsubsection{The algebra of distributions}
Let $k$ be a commutative unital ring. Suppose $X$ is an affine $k$-scheme, $x\in X(k)$ and set $$I_x=\{f\in k[X]\mid f(x)=0\}.$$ Then $k[X]=k1\oplus I_x$. A \emph{distribution of order $\leq n$} on $X$ with support at $x$ is a linear map $\mu:k[X]\to k$ with $\mu(I_x^{n+1})=0$. We set $\Dist_n(X,x)$ to be the set of all such; i.e. \[\Dist_n(X,x):=(k[X]/I_x^{n+1})^*\subseteq k[X]^*.\] Now set \[\Dist(X,x)=\{\mu\in k[X]^*\mid \exists n\in \mathbb{N}, \mu(I_x^{n+1})=0\}=\bigcup_{n\geq 0}\Dist_n(X,x).\] 
We say $X$ is \emph{infinitesimally flat} at $x$ if each $k[X]/I_x^{n+1}$ with $n\in \mathbb{N}$ is a projective $k$-module. Our main interest lies in case $X$ is $G$, and $x=1$ is the identity point. Then put $\Dist(G):=\Dist(G,1)$. Using the comultiplication we can make $\Dist(G)$ into an algebra; \emph{viz.}
\[\mu\nu:k[G]\stackrel{\Delta}{\longrightarrow}k[G]\otimes k[G]\stackrel{\mu\otimes\nu}{\longrightarrow}k\otimes k\stackrel{\sim}{\longrightarrow} k.\]
Hence we say $\Dist(G)$ is the \emph{algebra of distributions} on $G$. It is always true that a $G$-module inherits a (left) $\Dist(G)$-module structure, that a submodule of a $G$-module is a $\Dist(G)$-module and that homomorphisms of $G$-modules induce homomorphisms of $\Dist(G)$-modules, but not conversely. If $G$ is infinitesimally flat, noetherian and integral, then in fact a $\Dist(G)$-submodule $M'$ is a $G$-submodule whenever $M/M'$ is projective over $k$---\!\!\cite[I.7.15]{Jan03}; and \begin{equation}\Hom_G(M,M')=\Hom_{\Dist(G)}(M,M'),\label{eqn:disthoms}\end{equation} whenever $M'$ is projective over $k$---\!\!\cite[I.7.16]{Jan03}. [Recall $G$ is noetherian (resp.~integral) whenever $k[G]$ as a $k$-algebra is  noetherian (resp.~ an integral domain).]



\subsubsection{Socle, Radical, Rigidity}\label{sec:soc}
Let $R$ be a ring and $M$ an $R$-module of finite length; i.e.~a module with finitely many composition factors. Recall that the \emph{socle} $\Soc_R(M)$ of $M$ is defined to be the sum of its simple submodules. The \emph{socle series} (or \emph{Loewy series}) of $M$ is defined recursively by setting $\Soc^1M = \Soc_RM$
and letting $\Soc^iM$ be the submodule such that $\Soc^iM/\Soc^{i-1}M = \Soc_R(M/\Soc^{i-1}M)$. Dually, the \emph{radical} $\Rad_RM$ is the intersection of the maximal submodules of $M$. The radical series is defined recursively by letting $\Rad^1M = \Rad_RM$ and then $\Rad^iM:=\Rad_R(\Rad^{i-1}M)$. By convention, we set $\Rad^0(M) = M$. Taking $R$ as a module over itself, we have obviously $\Jac(R)=\Rad(R)$.

We say the socle series (resp.~radical series) has length $\ell:=\ell(M)$ if $\ell\in\mathbb N$ is minimal such that $\Soc^\ell(M)=M$ (resp.~$\Rad^\ell(M)=0$). 
They must have a common length $\ell$ called the \emph{Loewy length} and we always have an inclusion $\Rad^{\ell-i}(M) \subseteq \Soc^i(M)$ \cite[\S9.4]{ANT}.
 
If $M$ is indecomposable with Loewy length $\ell$---for example when it is the regular module for a local ring---then we say $M$ is \emph{rigid} if the radical and socle series coincide, that is $$
\Soc^iM = \Rad^{\ell-i}M \textrm{ for each } 0\leq i \leq \ell.$$ We say an $R$-module $M$ is \emph{rigid} if its indecomposable summands are so. 
If $R$ is a $k$-algebra for a field $k$, then $M$ is called \emph{geometrically rigid} if $M_{\bark}$ is a rigid $R_{\bark}$-module. It is called \emph{absolutely rigid} if $M_{E}$ is a rigid $R_E$ module for any field extension $E$ of $k$.


Following \cite[I.2.14, II.D.1]{Jan03}, we can replace $R$ with $G$ in all the above, where $G$ is an affine algebraic $k$-group and $k$ is a field. Then we get obvious notions of radical and socle series of finite-dimensional $G$-modules, and the idea of when one is rigid. 


The following collects some basic observations about rigidity, whose proofs are obvious:

\begin{lemma}\label{sumrigid}
Suppose $M$ is a finite-length $R$-module or $G$-module, and $M = U_1\oplus \cdots \oplus U_r$ is a decomposition of $M$ as a direct sum of submodules.
\begin{itemize}
\item[(i)] $\Soc^j(M) = \bigoplus_{i=1}^r \Soc^j(U_i) \textrm{  and } \Rad^j(M) = \bigoplus_{i=1}^r \Rad^j(U_i).$
\item[(ii)] The socle and radical series for $M$ coincide if and only if they coincide for each summand \emph{and} the summands all have a common Loewy length.
\item[(iii)] If the socle and radical series for $M$ coincide, then $M$ is rigid.
\end{itemize}
\end{lemma}

\subsubsection{On artinian algebras}\label{sec:artin} We need a little non-commutative algebra and our sources are \cite{Lam1} and \cite{Lam2}.
Let $k$ be a commutative artinian ring and $A$ a (possibly non-commutative) $k$-algebra which is finitely generated as a $k$-module (and hence artinian).
Let $\Jac(A)$ denote the Jacobson radical of $A$: the intersection of all maximal left ideals. It can be shown that $\Jac(A)$ is the annihilator of all simple left $R$-modules, from which it follows it is  a two-sided ideal of $A$---\hspace{1sp}\cite[\S4]{Lam1}.
Since $A$ is artinian, $\m_A:=\Jac(A)$ is the maximal nilpotent ideal of $A$ and the quotient $A/\m_A$ is the maximal semisimple quotient of $A$.
Further, an $A$-module $V$ is semisimple if and only if it is annihilated by $\m_A$---\hspace{1sp}\cite[Ex.~4.18]{Lam1}. 
So we see in this case that the terms of the socle series in a module $V$ are the annihilators of the powers of $\m_A$; \emph{viz.~}
\[\Soc^i(V)=\{v\in V\mid (\m_A)^i\cdot v=0\};\quad \Soc^i(A)=\Ann_A((\m_A)^i) ,\]
where we consider $A$ as a left regular $A$-module. If in addition $A$ is indecomposable \cite[\S22]{Lam1} and $n$ is the Loewy length of $A$, then we have
\begin{equation}\text{$A$ is rigid if and only if $(\m_A)^{n-i} = \Ann_A((\m_A)^i)$ for each $i$.}\label{rigidann}\end{equation}
If $A$ is not indecomposable, then $A$ is rigid if and only if the equality in \eqref{rigidann} holds for each of the blocks of $A$.
Note that, as in \Cref{sumrigid}, for a general artinian algebra $A$, if the equality in \eqref{rigidann} holds, then $A$ is rigid in such a situation;
the converse is not true in general (e.g., see \Cref{rem:notbottom}).


\begin{example}\label{rem:rigidegs}
We describe a non-rigid commutative artinian Gorenstein algebra. Let $k$ be any field and consider $A:=k[x,y]/(xy,x^3-y^2)$ as a module over itself. Then there is a unique maximal ideal $\m=\Jac(A)=(x,y)$.
So we see that $\Rad^i(A)$ coincides with the $i$th powers of $\m$, and these are $\m$, $\m^2=(x^2)$, and $\m^3=(x^3)$. 
Meanwhile, $\Soc^i(A)$ coincides with the annihilators $\Ann_A(\m^i)$; we have $\Soc(A)=(x^3)$, but $\Soc^2(A)=(x^2,y)$ and $\Soc^3(A)=(x)$. So even though $A$ is self-dual---i.e.~Gorenstein---$A$ is not rigid. Pictorially, its structure is as depicted in \cref{nonrigpic}, where submodules appear as a union of paths to the socle, which sits at the bottom. 

\begin{figure}\begin{center}\begin{tikzpicture}
  \matrix (m)
    [
      matrix of math nodes,
      nodes in empty cells,
      row sep    = 2em,
      column sep = 2em
    ]
    {
       & & \langle 1\rangle & \\
      \hspace{1em}\Rad(A) & \langle x\rangle & &  & \Soc^3(A)&[-2em]\\
       & & & \langle y \rangle & \Soc^2(A)\\
      \Rad^2(A)& \langle x^2\rangle & &  \\
        &\node(c){};& \Rad^3(A)=\langle x^3\rangle=\Soc(A) & & & \node(a) { }; \\
    };
    \draw (m-2-2) -- (m-1-3);
    \draw (m-2-2) -- (m-4-2);
    \draw (m-4-2) -- (m-5-3);
    \draw (m-1-3) -- (m-3-4);
    \draw (m-3-4) -- (m-5-3);
    \draw[{Bar}-{Bar}] (m-2-5) -| (a.center) |- (m-5-4);
    \draw[{Bar}-{Bar}] (m-3-5.south) |- (m-5-4.north east);
    \draw[{Bar}-{Bar}] (m-4-1.east) |- (c.center);
    \draw[{Bar}-{Bar}] (m-2-1.east) |- (c.north);
  
\end{tikzpicture}\end{center}\caption{A non-rigid algebra}\label{nonrigpic}\end{figure}
\end{example}

It was shown in \cite[\S70]{Mac} that one may characterise rigidity of a commutative Gorenstein  algebra of finite dimension over a field $k$ using its Hilbert function, so we explain this now. Take any filtration $\{M_i:=\mathscr{F}_i(M)\}$, i.e.~$M=M_0\supset M_1 \supset M_2 \supset \dots$, where the $M_i$ are all submodules of $M$. Then we may define the associated Hilbert function 
\[H_{\mathscr{F}}(M)(x):=H_{\mathscr{F},0}+H_{\mathscr{F},1}x+H_{\mathscr{F},1}x^2+\dots\] where $H_{\mathscr{F},i}=\dim_k(M_i/M_{i+1})$. 

Let $M=A$ and $\mathscr{F}$ the radical filtration $A_i:=(\m_A)^i$. 
Then Macaulay showed that the socle and radical series for $A$ coincide if and only the coefficients of $H_{\mathscr{F}}$ are symmetric about the middle. 
Rather similarly, the socle and radical series for $A$ coincide if and only if there is an isomorphism $A\cong\Gr_\m(A)$, where $\Gr_\m(A)$ denotes the associated graded algebra of $A$ arising from the radical filtration. There is a detailed study of the Hilbert functions that can arise in \cite{Iar94}.

\subsubsection{Base change}
Let $k$ be a field and $G$ an affine algebraic $k$-group; let $A$ be a finite-dimensional $k$-algebra. We record two results about the behaviour of $G$-modules and $A$-modules under base change. Let $V$ be an $A$-module. Then obviously for any field extension $E/k$, we have $V_E:=V\otimes_kE$ is an $A_E:=A\otimes_k E$-module in a natural way. If $W$ is another $A$-module then, analogously with \cref{eq:tensorhom}, the flatness of $E/k$ implies \[\Hom_A(V,W)_E\cong\Hom_{A_E}(V_E,W_E).\] Now we may state 

\begin{lemma}\label{lem:semisimpledown}
Suppose $E/k$ is a field extension and $V$ is a finite-dimensional $G$-module (resp. $A$-module).
If $V_{E}$ is semisimple as a $G_{E}$-module (resp.~ $A_E$-module), then $V$ is semisimple as a $G$-module (resp.~$A$-module).
\end{lemma}

\begin{proof}
We consider the case of $G$-modules; the proof for $A$-modules is identical. We prove the contrapositive.
Suppose that $V$ is not semisimple.
Then there is a non-split short exact sequence $0\to M \to V \to U \to 0$ with $U$ simple.
Since $U$ is simple, every element of $\Hom_G(U,V)$ must actually have image in the submodule $M$. 
Tensoring with $E$, we get from \eqref{eq:tensorhom} that $\Hom_G(U,V)_{E} = \Hom_{G_{E}}(U_{E},V_{E})$,
and since all the elements of $\Hom_G(U,V)$ have image in $M$, we have that all the elements of $\Hom_{G_{E}}(U_{E},V_{E})$
have image in $M_{E}$.
In other words, the short exact sequence $0\to M_{E} \to V_{E} \to U_{E} \to 0$ also has no splitting, so $V_{E}$ is not semisimple.
\end{proof}

The first part of the next result provides a partial converse to \cref{lem:semisimpledown}; it is well-known, but we include a proof for completeness.

\begin{lemma}\label{lem:sepext}
Suppose $E/k$ is a separable algebraic field extension, and $V$ is a $G$-module or an $A$-module.
\begin{itemize}
\item[(i)] $V$ is semisimple if and only if $V_{E}$ is semisimple.
\item[(ii)] If the Loewy length of $V$ is finite then $\Soc^i(V_{E})=(\Soc^i(V))_{E}$ and $\Rad^i(V_{E})=(\Rad^i(V))_{E}$.
\end{itemize}
\end{lemma}

\begin{proof} Again we give just the proof for $G$-modules. 

(i). By \cref{lem:semisimpledown}, we only need to prove the forward implication, for which it suffices to deal with the case that $V$ is simple.
Consider the separable closure $k_s$ of $k$, which contains $E$ since $E/k$ is algebraic.
The socle of $V_{k_s}$ is a non-trivial $G_{k_s}$-submodule of $V_{k_s}$, which is stable under the Galois group of $k_s/k$ and so
has a $k$-form: that is, there is a non-trivial $G$-submodule $U$ of $V$ with $U_{k_s} = \Soc_{G_{k_s}}(V_{k_s})$.
Since $V$ is simple, and $U$ is non-trivial, we have $U=V$.
Thus $V_{k_s}$ is semisimple.
Now the result follows from \cref{lem:semisimpledown} applied to $V_{E}$ and $V_{k_s}$.

(ii). By duality, it suffices to prove the first statement. And by induction, to prove the first statement, it suffices to prove $\Soc_G(V)_{E}=\Soc_{G_{E}}(V_{E})$. 
But that follows by the argument from (i).
\end{proof}

\subsubsection{Minimal fields of definition for Jacobson radicals} 
\begin{defn}Let $A$ be a finite-dimensional $k$-algebra, $K/k$ be a field extension and $M$ an $A_K$-module. An intermediate field $K/E/k$ is a \emph{field of definition} for $M$ if there is an $A_E$-module $N$ such that $N_K$ and $M$ are isomorphic $A_K$-modules. If $E$ is a field of definition for $M$ admitting no proper subfield of definition, then $E$ is a \emph{minimal field of definition} for $M$.

  Similarly, if $M$ is an $A$-module and $M'$ is an $A_K$-submodule of $M_K$, then $E$ is a field of definition for $M'$ if there is a submodule $N$ of $M_E$ such that $N_K=M'$, and it is minimal if $E$ admits no proper subfield of definition.
\end{defn}

\begin{remark}In general there is no guarantee of a minimal field of definition: in \cite{BR19} one can find an example of a two-dimensional module for the quaternion algebra $A:=\Q\{x,y\}/(x^2=y^2=-1,\ xy=yx)$ that is defined over the field $K=\Q(a,b)/(a^2+b^2+1)$, and whenever it is defined over some subfield $E$ of $K$, it is also defined over some proper subfield of $E$. \end{remark}

In the special case $A=k$, so that $M$ is a $K$-vector space, one sees $M\cong N_K$ if $N$ is a $k$-vector space with a basis of the same cardinality as that of $M$ so that its minimal field of definition exists and is $k$. When $V$ is any $k$-vector space and $W$ a $K$-subspace for some field extension $K/k$, it is explained in \cite[Rk.~1.1.7]{CGP15} how to construct the unique minimal field of definition $E$ for $W$: one takes a basis $\{e_i\}_{i\in I}$ of $V$, with a subset $\{e_j\}_{j\in J}$ that maps to a basis $\{\overline{e_j}\}_{j\in J}$ of $(V\otimes K)/W$; then one takes for $E$ the subfield of $K$ spanned by the coefficients of the remaining $\{\overline{e_i}\}_{i\in I\setminus J}$ when expressed as linear combinations of the $\overline{e_j}$. At least for an artinian algebra $A$ this construction is compatible with the formation of the Jacobson radical of $A_\bark$ as we describe below. 


\begin{lemma}\label{lem:mfod}
 Let $k$ be a field and let $A$ be a finite-dimensional algebra over $k$.
 Let $\mathcal{J}:=\Jac(A_{\bark})$ denote the Jacobson radical of the base change $A_{\bark} = A\otimes_k \bark$.
 \begin{itemize}
  \item[(i)] If $A=K$ is a purely inseparable field extension of $k$, then $\mathcal{J}$ has minimal field of definition $K$. 
 \item[(ii)] More generally, the minimal field of definition of $\mathcal{J}$ is a finite purely inseparable extension $K/k$.\end{itemize}
 Let $J\subseteq A_{K}$ be such that $\mathcal{J} = J_\bark$.
\begin{itemize}\item[(iii)]  For any extension $E/K$ we have $\Jac(A_E)=J_E$.
\item[(iv)] The module $A_K$ is rigid if and only if $A_E$ is rigid for some field extension $E/K$ if and only if $A_E$ is rigid for all field extensions $E/K$.
\end{itemize}
\end{lemma}

\begin{proof}
  (i). Because $K$ is purely inseparable, for any algebraic extension $E/k$ the algebra $K\otimes_k E$ is local, since there is precisely one embedding $K\hookrightarrow \overline{E}=\bark$ and so there is just one possible quotient field: the compositum of $K$ and $E$. 
  Seen as an $E$-algebra this quotient field is $1$-dimensional precisely when $K\subseteq E$, hence $\dim \Jac(K\otimes_k \bark)=[K:k]-1$ and $K$ is the required minimal field of definition.
  
(ii). Since $A$ is finite-dimensional, by the Artin--Wedderburn theorem we have $$A_{\text{ss}}:=A_{k_s}/\Jac(A_{k_s})\cong\prod_{i=1}^r\Mat_{n_i}(k_i)$$ for $n_i\in\N$ and each $k_i$ a finite purely inseparable field extension of $k_s$. By (i) we have that $k_i$ is the minimal field of definition for $\Jac(\Mat_{n_i}(k_i)_\bark)=\Mat_{n_i}(\Jac((k_i)_\bark))$ and so the compositum $L$ of the $k_i$ is the finite purely inseparable extension of $k_s$ which is the minimal field of definition for $\Jac(A_{\text{ss}})_\bark$. Indeed we have $(A_{\text{ss}})_L/\Jac((A_{\text{ss}})_L)\cong\prod\Mat_{n_i}(L)$, which is an absolutely semisimple $L$-algebra expressed as a quotient of $A_L$ by a nilpotent ideal. Hence $\mathcal{J}$ is defined over $L$. Conversely, if $E$ does not contain some $k_i$, then there is some quotient algebra $\Mat_{n_i}(k_i)_E$ of $A_E$ whose Jacobson radical is $\Mat_{n_i}(\Jac(k_i\otimes_k E))$; by (i), $\Jac(k_i\otimes_k E)_\bark$ is a strict subalgebra of $\Jac(k_i\otimes_k \bark)$ and so $\mathcal{J}$ is not defined over $E$.

Since $\bark\cong \kp\otimes k_s$ where $\kp=k^{p^{-\infty}}$, there is a unique $k$-descent of $L$; i.e.~there is a purely inseparable field extension $K/k$ such that $K\otimes_k k_s\cong L$. Since $\Jac(A_L)$ is a characteristic ideal, it is stable under the absolute Galois group $\Gal(k_s/k)\cong\Gal(L/K)$, and so by Galois descent has $K$ as its minimal field of definition. 

(iii). We already observed that $A_L/J_L=(A_K/J_K)_L$ is absolutely semisimple, and so $A_K/J_K$ is absolutely semisimple also. 

(iv). Since any artinian algebra is a direct product of local artinian algebras whose factors are the indecomposable summands of $A$ as a left $A$-module, we may assume it is local, since rigidity of a module is predicated on its indecomposable summands. Let $E/K$ be any field extension.
Since $A_E$ is artinian, the radical series for $A_E$ is formed by taking the powers of $\Jac(A_E)$, and the socle series by taking the annihilators of those powers.
Part (iii) says that $\Jac(A_E) = J_E$, so the powers of $\Jac(A_E)$ are the base changes to $E$ of the powers of $J$.
Then the annihilators of the powers of $\Jac(A_E)$ are the base changes to $E$ of the annihilators of the powers of $J$ as well: certainly we have an inclusion $\Ann(J^i)_E \subseteq \Ann(J_E^i)$, and then a consideration of dimension gives equality.  
Hence, if the socle and radical series coincide over $E$, they already coincide over $K$, and vice versa.
\end{proof}


\begin{remark}\label{mfodsep}
Let us underline the aspect of the construction of $K$ in the proof which shows that the minimal field of definition of $\mathcal{J}$ commutes with separable extensions. Suppose $J\subseteq A_K$ such that $J_\bark=\mathcal{J}$.
If $E/k$ is some separable extension,  $E':=K\otimes_k E$ is a field since $K/k$ is purely inseparable and $E/k$ is separable, and $\mathcal{J}$ is $E'$-defined via the ideal $J_{E'}$ of $A_{E'}$.
Thus the minimal field of definition over $E$ of $\mathcal{J}$ is contained in $E'$.
On the other hand, if $L/E$ is any extension over which $\mathcal{J}$ is defined, then $L/k$ is also an extension over which $\mathcal{J}$ is defined, so $L$ contains $K$.
Hence $L$ contains $E'$ and $E'$ is the minimal field of definition of $\mathcal{J}$ over $E$. 
\end{remark}


\subsection{Two Morita equivalences}\label{sec:morita}
Let $k$ be a commutative ring and suppose $G$ is flat,  infinitesimally flat, noetherian, and integral. Let $M$ be a $G$-module which is projective as a $k$-module. 

Let $\phi:\Dist(G)\to T:=\End_k(M)$ be the action map, and $R$ the quotient of $\Dist(G)$ by $\ker \phi$, identifying $R$ with a subalgebra of $\End_k(M)$, such that $M$ is a faithful  $R$-module. We have a short exact sequence
\[0\to \ker\phi\to\Dist(G)\to R\to 0,\] 
and if $E$ is a flat $k$-algebra, then the sequence
\[0\to (\ker\phi)_E\to\Dist(G)_E\to R_E\to 0,\] 
is exact. Now \cite[I.7.4(1)]{Jan03} gives $\Dist(G_E)\cong \Dist(G)_E$, hence the action of $\Dist(G_E)$ factors through $\phi_E$---indeed, as $E$ is flat, we have $(\ker\phi)_E=\ker(\phi_E)$; so $R_E$ is faithful on $M_E$. 
Now we get from \cref{eqn:disthoms} that the $G_E$-submodules of $M_E$ are precisely the $R_E$-submodules of $M_E$.

We apply the above to the case where $M$ is a simple $G$-module which is finitely generated as a $k$-module (for example, when $k[G]$ is projective or $k$ is noetherian). In that case, $M^*$ is finitely generated by the dual basis lemma \cite[2.11]{Lam2}; moreover, we also have $\End_k(M)\cong M^*\otimes_k M$ \cite[Ex.~\S2.20]{Lam2} and so $\End_k(M)$ is finitely generated (and projective as a $k$-module). Now Schur's lemma tells us that $\End_G(M)\cong \End_{\Dist(G)}(M)\cong\End_R(M)=:D$ is a division algebra over $k$, and it is finitely-generated as a $k$-module, hence artinian. As $R$ is left primitive (i.e.~has a faithful left module), it must be simple---\!\!\cite[11.7]{Lam1}. 

An $R$-module $P$ is said to be a \emph{left generator} for $R$ if $\Hom_R(P,?)$ is a faithful functor from the category of left $R$-modules to the category of abelian groups---\!\!\cite[18.7]{Lam2}. If in addition $P$ is finitely generated and projective it is called a \emph{progenerator} for $R$. 
As $R$ is simple, the category of left $R$-modules is semisimple and so any nonzero module is a generator; 
thus $M$ is a progenerator of $R$. Since $D$ is the centraliser of $R$ in $T$, we get that $M$ is a right $D$-module and $M$ is an $(R,D)$-bimodule that is faithfully balanced \cite[18.21]{Lam2}; this is to say the maps $R\to \End_D(M)$ and $D\to \End_R(M)$ are both ring isomorphisms---so $R$ is also the centraliser in $T$ of $D$. In particular, $R$ and $D$ are Morita equivalent---\!\!\cite[18.33]{Lam2}.

Now, under flat base change through $k\to E$,  we have  that $D_E\cong \End_{R_E}(M_E)$.
We also have that $M_E$ is a progenerator for $R_E$: it is finitely-generated projective since $M$ is and $E$ is flat; and it is a generator
by the characterisation in \cite[18.8(3)]{Lam2} applied to $M$ (resp. $M_E$) and $R$ (resp. $R_E$). 
Thus $M_E$ is a faithfully balanced $(R_E,D_E)$-bimodule, so  $R_E$ and $D_E$ are again Morita equivalent and
 \cite[18.44]{Lam2} furnishes us with:
\begin{prop}\label{prop:REDE}
The  lattice of (left) $R_E$-submodules of $M_E$ is isomorphic to the lattice of left ideals of $D_E$.
\end{prop}

We can push this analysis one step further.
Let $Z :=Z(D)$ be the centre of the division ring $D$.
Considering $D$ as a left $D$-module, multiplication on the right by elements of $D$ gives an identification $\End_D(D) \cong D$, 
and the centraliser of $D$ is just the centre $Z$.
It is clear that we get another Morita equivalence with $D$ itself as the progenerator this time, and we deduce:
\begin{cor}\label{cor:DEZE}
The  lattice of left ideals of $D_E$ is isomorphic to the lattice of ideals of $Z_E$.
\end{cor}

\begin{center}\framebox{For the rest of the paper, $k$ will denote a field.}\end{center}

In particular, since $k$ is a field, then $G$ satisfies the hypotheses at the beginning of the section; so the conclusions of \cref{prop:REDE} and \cref{cor:DEZE} both hold. The following result  also applies to algebras over a field and enables the use of the Conrad--Prasad classification.

\begin{lemma}\label{directproduct}Suppose $A$ and $B$ are two  associative unital $k$-algebras and put $C:=A\otimes_k B$, where $A$ and $B$ are considered commuting subalgebras of $C$ via $A\cong A\otimes 1$ and $B\cong 1\otimes B$. 
  Let $W$ be a finite-dimensional simple $C$-module. Then the following hold:
 \begin{enumerate}
   \item $W|_{A}$ is an  isotypic semisimple $A$-module, say $U^r$ for a simple $A$-module $U$, and $r\in\NN$; similarly let $W|_{B}\cong V^s$. There is a surjective $C$-module homomorphism $\psi:U\otimes V\to W$.\label{ds3}
 \end{enumerate}
   Suppose $D:=\End_A(U)$, $E:=\End_B(V)$ and $F:=\End_C(W)$. Then:
   \begin{enumerate}
   \setcounter{enumi}{1}
   \item $\End_C(U\otimes V)\cong D\otimes E$;\label{ds1}
   \item $Z(D\otimes E)\cong Z(D)\otimes Z(E)$;\label{ds2}
   \item $C/\Ann_C(W)$ is a simple $k$-subalgebra of $\End_k(W)$ generated by the images of $A\otimes 1$ and $1\otimes B$; \label{ds4}
   \item if $k=k_s$, then $Z(F)$ is the compositum of the purely inseparable field extensions $Z(D)$ and $Z(E)$;\label{ds5}
   \item for arbitrary $k$, the following fields coincide\label{ds6}:
   \begin{enumerate}\item the minimal fields of definition of $\Jac(F_\bark)$ and $\Jac(Z(F)_\bark)$; 
     \item the compositum of the minimal fields of definition of $\Jac(D_\bark)$ and $\Jac(E_\bark)$;
     \item the compositum of the minimal fields of definition of $\Jac(Z(D_\bark))$ and $\Jac(Z(E_\bark))$.\label{dstest}
   \end{enumerate}
 \end{enumerate}
\end{lemma}
\begin{proof} 
 \cref{ds3} is \cite[VIII, ~\S12.1, Prop.~2]{BourbAlg8}, but we nutshell the details. One takes a simple $A$-submodule $U$ of $W$ and considers $X:=\Hom_A(U,W)$, which becomes a $B$-module via $(b\circ\phi)(u)=b(\phi(u))$. Let $\varphi:V\to X$ be the embedding into $X$ of some simple $B$-submodule. Then define $\psi:U\otimes V\to W$ by $\psi(u\otimes v)\mapsto \psi(v)(u)$. Then one checks $\psi$ is a non-zero $C$-module map, and the isotypicity of $W|_{A}$ follows from that of $(U\otimes V)|_A\cong U^{\dim V}$.
 
Now \cref{ds1} and \cref{ds2} are \cite[VIII, \S12.5, Prop.~5(a)] {BourbAlg8}. By Morita theory---see \cref{prop:REDE}---the fact that $W$ is simple implies that $C/\Ann_C(W)$ is simple, giving \cref{ds4}. 

Take $k=k_s$. Precisely the same argument applied with $(D,E,F,Z(D),Z(E),Z(F))$ replacing $(A,B,C,D,E,F)$ yields that $Z(F)$ is a simple quotient of $Z(D)\otimes_k Z(E)$. Since $k=k_s$, $Z(E)\otimes_k Z(F)$ is a tensor product of purely inseparable extension fields and so it is local.
Thus $Z(F)$ must identify with its quotient field, which is of course the compositum of $Z(E)$ and $Z(F)$ as claimed in \cref{ds5}.

Since the minimal fields of definition of the Jacobson radicals commute with separable extension, we may base change to $k_s$, whereupon $W$ is still semisimple, and argue with each simple submodule individually. This gives us a collection of division rings $(D_i,E_i,F_i)$ and the minimal fields of definition of the Jacobson radical of $D_{k_s}:=\prod D_i$ is the compositum of all the $Z(E_i)$ and $Z(F_i)$---\cref{lem:mfod}(i). This proves \cref{ds6}.
\end{proof}

Suppose $H_1$ and $H_2$ are smooth connected algebraic $k$-groups and let $J:=H_1\times H_2$. Applying the lemma to the case $A:=\Dist(H_1)$, $B:=\Dist(H_2)$, and $C:=\Dist(J)\cong\Dist(H_1)\otimes\Dist(H_2)$---see \cite[I.7.9(3)]{Jan03}---we get the following.

\begin{cor}\label{redtofactors}
  Suppose $V$ is a simple $J$-module. 
Then the restriction of $V$ to $H_i$ decomposes as a direct sum of copies of a single simple module $W_i$ for $i=1,2$.
Denote by 
$$D:=\End_{J}(V),\ D_1:=\End_{H_1}(W_1),\ D_2:=\End_{H_2}(W_2),$$ the respective division rings. 
Then if $K$, $K_1$, $K_2$ are the respective minimal fields of definition of their geometric Jacobson radicals (or equivalently, the geometric Jacobson radicals of the centres $Z(D)$, $Z(D_1)$ and $Z(D_2)$), then $K$ is the compositum of $K_1$ and $K_2$.
\end{cor}


By \cref{eq:tensorhom} and the earlier remarks in this section, $\End_{G_E}(M_E)\cong \End_{R_E}(M_E)\cong D_E$ and the lattice of $G_E$-submodules of $M_E$ is the same as the lattice of $R_E$-submodules,  thus isomorphic to that of $D_E$ as a left $D_E$-module, and hence to that of $Z_E$ as a $Z_E$-module.  
From all of this, we get
\begin{cor}\label{cor:GEDE}
Let $M$ be a simple $G$-module, and $E/k$ an extension of fields.
The $G_E$-module $M_E$ is rigid if and only if the commutative algebra $Z_E$ is rigid (as a module for itself). 
\end{cor}

\begin{proof}
The hypotheses on $k$ and $G$ ensure that all the preceding results in this section hold.
But then we see that the socle and radical series for $M_E$ as an $R_E$-module correspond to the socle and radical series for $Z_E$, by \Cref{prop:REDE} and \Cref{cor:DEZE}
applied to this case.
\end{proof}

Thus, to complete the proof of \cref{thm:mainone} for connected $G$, we need to show the existence of a purely inseparable extension $K/k$ such that $Z_E$ is rigid for every extension $E/K$. 
This is achieved in the next section, where we show that we can take $K$ to be the minimal field of definition of $\Jac(Z_\bark)$.

Finally in this section we give some further generalities about semisimple $k$-algebras, which will help us when we are describing
the endomorphism rings of simple modules in Part \ref{partii}.

\begin{defn}
  Let $A$ be a finite-dimensional semisimple $k$-algebra over a field of characteristic $p$. 
  We say $A$ is \emph{$p$-split} if $A$ is $k$-isomorphic to a direct product of matrix rings over purely inseparable field extensions of $k$.

  If $A$ is a $k$-algebra and $E$ a field extension of $k$ such that $A_E:=A\otimes_kE$ is $p$-split as an $E$-algebra, then we say that $A$ is \emph{$p$-splittable}. 
  If $[E:k]$ minimal such that $A_E$ is $p$-split, then we say $E$ is a \emph{$p$-splitting field for $A$}.

  If $A$ is a simple $k$-algebra whose centre $Z(A)$ is a purely inseparable extension of $k$, then we say that $A$ is a \emph{$p$-central simple algebra} (or \emph{$p$-CSA}). 
  If $A$ is also a division algebra then we say $A$ is a \emph{$p$-division algebra}.
\end{defn}

We aim to show:

\begin{theorem}
 Let $A$ be a finite-dimensional semisimple $k$-algebra and $Z:=Z(A)$ its centre. Then: 
\begin{itemize}
\item[(i)] the algebra $A$ is $p$-split by a field $\mathscr{E}$ that is finite and Galois over $k$;
\item[(ii)] if $A$ is simple then any $p$-splitting field $E$ for $A$ contains the normal closure $F$ of the separable part $Z_\sep$ of the field extension $Z/k$. If $\ell:=[Z_\sep:k]$ then $$Z(A_F)\cong \underbrace{F\times\dots\times F}_{\ell-\text{times}},\qquad\text{and}\qquad A_F\cong A_1\times\dots\times A_\ell$$ as a direct product of $p$-CSAs. 
\end{itemize}
\end{theorem}

\begin{proof}
  The existence of a field extension which $p$-splits $A$ is obvious: taking $E=k_s$ one gets that $A_E$ is still semisimple and is therefore still a product of matrix rings over division $k_s$-algebras by Artin--Wedderburn. But the Brauer group of $k_s$ is trivial; so all these division algebras are purely inseparable field extensions of $k_s$.
    
We first prove the rest of the result for a division $k$-algebra $D$; the general result follows quickly from Artin-Wedderburn.
Suppose $Z=Z(D)$, let $Z_\sep$ be the separable part of this field extension of $k$, and
let $E/k$ be some field extension that $p$-splits $D$; so $D_E:=D\otimes_k E$ is a direct product of matrix algebras over a set  of finite purely inseparable field extensions of $E$; and its centre $Z(D_E)$ isomorphic to the direct product of these fields. Take any $x\in Z_\sep$ with minimal polynomial $f$ over $k$. Then $k[x]/f$ identifies with a subfield of $Z_\sep$ and therefore $(k[x]/f)\otimes_k E\cong E[x]/f$ identifies as an $E$-subalgebra of $Z_E$. 
  But since $f$ is separable, $E[x]/f$ is also isomorphic to a direct product of separable extensions of $E$, and since $Z_E$ is a product of purely \emph{in}separable extensions of $E$ we see that $E$ contains all the roots of $f$.
  Thus $E$ contains the normal closure $F$ of $Z_\sep$ as claimed. This proves the first part of (ii).

  For the second, note that we have 
  $$
  Z_F= Z\otimes_k F \cong Z\otimes_{Z_\sep} (Z_\sep\otimes_k F).
  $$ 
  Writing $Z_\sep = k[x]/f$ for some separable polynomial $f$ and using the fact that the roots of $f$ all lie in $F$ by the previous paragraph, we see that $(Z_\sep\otimes_k F)$ is isomorphic to a product of $\ell=[Z_\sep:k]$ copies of $F$, conjugate under the Galois group of $F/k$. 
  Since $Z/Z_\sep$ is purely inseparable and $F/Z_\sep$ is separable, $Z$ and $F$ are linearly disjoint; so $Z_F$ is isomorphic to $\ell$ copies of the compositum $ZF:=(Z\otimes_{Z_\sep}F)$, and these fields are still conjugate under $\Gal(F/k)$. 
  Now $D_F$ is semisimple because $F/k$ was separable; so $$D_F\cong A_1\times \dots\times A_\ell$$ where the $A_i$ are $p$-CSAs that are all conjugate under $\Gal(F/k)$: 
 this decomposition follows from the decomposition of $1\in Z_F$ into central primitive idempotents corresponding to the $\ell$ Galois-conjugate factor fields of $Z_F$, see \cite[7.22]{Lam2}. This proves (ii) in our special case; and the general case of (ii) is also clear as $A\cong \Mat_r(D)$ for some division $k$-algebra $D$.
  
  By Galois conjugacy, there are integers $n$ and $d$ such that any of the $p$-CSA factors of $D_F$, say $A_1$, has $A_1\cong \Mat_{n}(D_1)$ and $D_1$ is some division $F$-algebra of degree $d$---i.e.~$\dim_F D_1=d^2$---and whose centre $Z(D_1)$ is isomorphic to $ZF$.
  Now, $A_1$ and $D_1$ are both central simple as $ZF$-algebras. Therefore by \cite[15.12]{Lam1} they are both split by a maximal subfield of $D_1$ that is separable over $ZF$, say $K_1$, where $[K_1:ZF]=d$. 
  Since $ZF/F$ is purely inseparable, $K_1/ZF$ is the separable part of the extension $K_1/F$, hence $K_1/F$ descends to some separable $K_1'/F$.
  Then $K_1\cong ZF\otimes_F K_1'$, and $K_1/K_1'$ is purely inseparable.
  We conclude that the semisimple algebras $A_1\otimes_F K_1'$ and $D_1\otimes_F K_1'$ are both $p$-split as $K_1'$-algebras, since
  $$
  D_1\otimes_{ZF} K_1' \cong D_1\otimes_{ZF} (ZF\otimes_F K_1') \cong D_1\otimes_F K_1'.
  $$
  Taking the normal closure of $K_1'$ in $k_s$ gives a Galois extension $E/k$. If $\sigma\in \Gal(F/k)$ is such that $\sigma(A_1)=A_i$ then extend it to $\hat\sigma\in \Gal(E/k)$, and put $K_i':=\sigma(K_1')$. Since $(A_1)_{K_1'}$ is $p$-split, so is $(A_i)_{K_i'}$ and we get that all $A_i$ are $p$-split by $E$. Evidently, Artin--Wedderburn  implies (ii) for the more general simple $k$-algebra $A$.
  
  It remains to see the existence of the Galois extension in part (i), when $A$ is semisimple. Apply the above to each simple factor of $A$ to get a collection of Galois extensions of $k$ and take their compositum $\mathscr{E}$. We have that $A_\mathscr{E}$ is semisimple and each of its simple factors is $p$-split; hence $A_\mathscr{E}$ is $p$-split too.
\end{proof}

\begin{remark}\label{rem:noncan}While the field $F$ in part (ii) of the theorem is uniquely determined, the field $E$ is not. For example when $A=\Q\oplus i\Q\oplus j\Q\oplus k\Q$ is the $\mathbb{Q}$-division algebra of Hamiltonian rational quaternions, then it is ($p$-)split by many quadratic field extensions. Moreover,  following Amitsur's original construction, there are by now many examples of finite-dimensional central $k$-division algebras $D$ which are \emph{non-crossed}; that is to say that there is no maximal subfield of $D$ that is Galois. From this it also follows that $D$ cannot be constructed from a cocycle using Noether's method---see the introduction of \cite{Han04} for an overview. In any case, the point to make here is that the field $\mathscr{E}$ is highly non-canonical in general.

  In the scenario of interest to us---namely when $D=\End_G(V)$ for $V$ a finite-dimensional $G$-module---we will show by contrast that there is a \emph{unique} $p$-splitting field $\mathscr{E}$, and that is Galois over $k$; see \cref{thm:torussplit}.
\end{remark}



\subsection{On the rigidity of finite-dimensional algebras}\label{sec:rigidity}

\begin{theorem}\label{thm:Drigid}
Suppose that $R$ is a finite-dimensional simple $k$-algebra (though not necessarily central simple).
Let $k'/k$ be the minimal field of definition of $\Jac(R_{\bark})$. Then $R_{k'}$ is absolutely rigid. 
\end{theorem}

Note that the proof of this theorem immediately reduces to the case that $R$ is a field:
since $R$ is simple, the Artin-Wedderburn Theorem says that $R\cong M_n(D)$ for some $D$, where $D$ is a finite-dimensional division $k$-algebra.
For any field extension $E/k$, we have $R_E \cong M_n(D_E)$, and the ideal structure of $R_E$ is therefore identical with that of $D_E$---\cite[3.1]{Lam1}.
But now we may replace $D$ with its centre $Z$, as in \Cref{cor:DEZE}.
Since $R$ is finite-dimensional, $Z/k$ is some finite extension of $k$.
We can also see that the minimal fields of definition of $\Jac(R_\bark)$, $\Jac(D_\bark)$ and $\Jac(Z_\bark)$ all coincide.
 



\subsubsection{Generalities on tensor products of fields}\label{sec:tensor}
The literature already contains a number of results about the tensor product $A:=K\otimes_k E$ of general field extensions $K/k$ and $E/k$. 
For example, it is a result of Grothendieck, in generalised form by Sharp \cite{Sha77} that $\dim(K\otimes_k E)=\min(\mathrm{tr.deg}(K/k),\mathrm{tr.deg}(E/k))$. 
Furthermore Grothendieck proved in \cite[Lem.~6.7.1.1]{EGAIV65} that if one of the extensions is finite, then $A$ is Cohen-Macauley. This has been generalised in at least two directions: in \cite[Lem.~2.2]{BK02} weakening the hypotheses to demanding $A$ be noetherian; and in \cite[I.Thm.~2]{watanabe} strengthening the conclusion to saying $A$ is  Gorenstein. (Recall that a zero-dimensional commutative local noetherian ring $A$ is \emph{Gorenstein} if one of the following equivalent conditions holds \cite[Prop. 21.5]{Eisenbud}: $A$ has a simple socle as a left $A$-module; $A$ is self-injective. Or if $A$ is a $k$-algebra then equivalently $A$ is self-dual over $A$, that is, $A\cong \Hom_k(A,k)$.)

A \emph{simply truncated polynomial algebra (STP algebra)} over a field $k$ is an algebra of the form
$$
A=k[x_1,\ldots,x_n]/(x_1^{a_1},\ldots,x_n^{a_n})
$$
with $a_1\geq a_2 \geq \cdots\geq a_n$ \cite[Ch. 1]{rasala}.
Let $x_i$ denote the image of $X_i$ in $A$.
Then $A$ is a local ring, with maximal ideal $\m$ generated by the $x_i$.
It is clear that if $A$ is an STP algebra then, for any field extension $E/k$, $A_E = A\otimes_k E$ is an STP algebra over $E$.

\begin{lemma}\label{rigidSTP}
An STP algebra $A$ is a rigid local Gorenstein algebra, and hence has a symmetric Hilbert function.
\end{lemma}

\begin{proof}
Since $A$ is local, it is indecomposable and thus we need to check \eqref{rigidann}.
For a tuple $\beta = (b_1,\ldots,b_n)$ of non-negative integers, we let $x^\beta := \prod_{i=1}^n x_i^{b_i}$,
and for another such tuple $\beta' = (b_1',\ldots,b_n')$ we say $\beta\leq \beta'$ when $b_i\leq b_i'$ for every $i$.  
Note that for any $\beta = (b_1,\ldots,b_n)$, we have $x^\beta = 0$ if and only if there exists some $i$ with $b_i\geq a_i$.
Let $\gamma = (a_1-1,\ldots,a_n-1)$, and let $n = \left(\sum_{j=1}^n (a_j-1)\right) + 1$.
Since $x^\gamma\in\m^{n-1}$ we have $\m^n = \{0\}$ but $\m^{n-1}\neq \{0\}$.
We need to show that $\Ann(\m^i) = \m^{n-i}$ for each $1\leq i\leq n-1$.
It is clear that for any $1\leq i\leq n-1$, $\m^{n-i}$ annihilates $\m^i$.
On the other hand, given $1\leq i\leq n-1$, the ideal $\m^{n-i-1}$ is generated by the $x^\alpha$ with $\alpha=(\alpha_1,\ldots,\alpha_n) \leq \gamma$ and $\sum_{j=1}^n \alpha_j = n-i-1$.
Given any such $\alpha$, let $\beta = \gamma-\alpha$.
Then $x^\alpha x^\beta = x^\gamma \neq 0$, and $\sum_{j=1}^n \beta_j = i$, so $x^\alpha \not\in \Ann(\m^i)$.
Thus $\Ann(\m^i) = \m^{n-i}$, and we see $A$ is rigid.

The final power $\m^{n-1} = \Ann(\m) = \Soc(A)$ is generated by the element $x^\gamma$, and so is simple as an $A$-module, so $A$ is Gorenstein.
The symmetry of the Hilbert function of $A$ follows from \cite[\S70]{Mac}, as explained in \Cref{rem:rigidegs}.
\end{proof}

Now suppose $Z/k$ is some finite extension, and let $M\subseteq Z$ be the separable part of the extension.
From \cite[Ch.~2, Thm.~6]{rasala} we learn that there is some finite (normal) extension $E/Z$ such that $Z\otimes_k E$ is a sum of STP $E$-algebras, and the summands are all isomorphic (via Galois automorphisms) to $Z\otimes_M E$.
In particular, $Z\otimes_k \bark$ is a sum of STP $\bark$-algebras, and all the summands are isomorphic as rings.

\begin{prop}\label{stp}
Let $k'$ be the minimal field of definition of $\Jac(Z_\bark)$.
Then $Z_{k'} := Z\otimes_k k'$ is absolutely rigid.
Thus for all extensions $E/k'$, the Hilbert functions $H(Z_E)$ coincide, and are symmetric about the middle degree term.
\end{prop}

\begin{proof}
By the results of \cite{rasala} above, $Z_\bark$ is a sum of isomorphic copies of the STP $\bark$-algebra $Z\otimes_M \bark$;
these will be the blocks of $Z_\bark$.
Since by \cref{rigidSTP} each of these is rigid, so is $Z_\bark$. 
In fact, since the blocks are all isomorphic as rings, we are in the sitiuation where the socle/radical series of $Z_\bark$ is given by powers of its Jacobson radical.
In any case, we can now apply 
\Cref{lem:mfod}(iii) with $A=Z$ to deduce rigidity of $Z_E$ for all extensions $E/k'$.

The final statement about the Hilbert functions follows since the socle and radical series for any $Z_E$ are just the base changes of those for $Z_{k'}$, and 
so all the Hilbert series coincide.
The dimensions of quotients for that series can therefore be calculated over $\bark$, and since $Z_\bark$ is the sum of isomorphic copies of an STP, \Cref{rigidSTP} gives us the symmetry result.
\end{proof}

The proposition above completes the proof of \cref{thm:Drigid}, recalling the observations at the start of this section. It also completes the proof of \Cref{thm:mainone} for connected $G$ using the Morita equivalences of \cref{sec:morita}:
if $V$ is a simple $G$-module, then $V_{k_V}$ is absolutely rigid,
where $k_V$ is the minimal field of definition of $\Jac(\End_G(V)_\bark)$.

\subsubsection{Example: a non-rigid tensor product of fields}
We give an example of a tensor product of two finite purely inseparable field extensions $K/k$, $E/k$ whose regular module is not rigid,
which shows that \cref{thm:Drigid} does not hold in full generality without the extra hypothesis involving the minimal field of definition $k'$.

From this, \cref{rem:finally} implies the existence of a simple module for a pseudo-split pseudo-algebraic group which is not absolutely rigid.

\begin{example}\label{ex:finally}
Let $k=\F_2(a,b,c,d)$ where $a,b,c,d$ are algebraically independent transcendental elements.
Consider the following purely inseparable extensions $K/k$ and $E/k$
\begin{align*}
  K&:=k(\underbrace{a^{1/16}+b^{1/4}}_\beta,\underbrace{a^{1/8}+c^{1/4}}_\gamma,\underbrace{a^{3/16}+d^{1/4}}_\delta) = k(\beta,\gamma,\delta)\\
  E&:=k(\underbrace{a^{1/16}}_{f_1},\underbrace{a^{3/16}+a^{1/8}b^{1/4}+a^{1/16}c^{1/4}+d^{1/4}}_{f_2}) = k(f_1,f_2),
\end{align*}
and let $A = K\otimes_k E$ with maximal ideal $\m_A$.

Note that $\gamma^4 = a^{1/2} + c = \beta^8+b^2+c$, and $\delta^4 = a^{1/4} + d = \beta^4 + b + d$, so $[K:k] = 16\times 4\times 4 = 2^8$.
We also have $[E:k]= 16\times 4 = 2^6$, and hence $\dim_k(A) = 2^{14} = 16,384$.
The compositum of $K$ and $E$ is the field 
$$
KE = k(a^{1/16},b^{1/4},c^{1/4},d^{1/4}) \cong A/\m_A
$$ 
of degree $2^{10}$ over $k$, so the maximal ideal $\m$ has dimension $2^{14} - 2^{10} = 2^{10}\times 15$ over $k$.
Thus $A$ has $16$ composition factors isomorphic to $KE$, with $15$ of them coming from $\m$.

Viewing $A$ as a $K$-algebra through multiplication in the first factor, $A$ has a $K$-basis consisting of
the $2^6$ elements $1\otimes f_1^if_2^j$ where $0\leq i\leq 15$ and $0\leq j\leq 3$.
An element $x = \sum_{i,j} e_{ij}\otimes f_1^if_2^j$ lies in $\m_A$ if and only if $\sum_{i,j} e_{ij}f_1^if_2^j = 0$ in $KE$.
In other words, the elements of $\m_A$ correspond to $K$-linear dependences between the $f_1^if_2^j$.
Noting that $f_1^4=a^{1/4}=\beta^4+b\in K$, and $f_2 = \beta f_1^2 + \gamma f_1 + \delta  \in K(f_1)$, it is not hard to show that the elements
$$
m_1:=1\otimes f_1^4+f_1^4\otimes 1 \qquad \textrm{ and } \qquad m_2:=1\otimes f_2 + \beta\otimes f_1^2+\gamma\otimes f_1+\delta\otimes 1
$$
generate $\m_A$.

Some straightforward calculations show that the first power of $m_1$ that is $0$ is $m_1^4=0$, the first power of $m_2$ that is $0$ is $m_2^8= 0$,
and $m_2^4 = m_1^3$.
Thus we have the following:
\begin{center}\begin{minipage}[t]{0.4\textwidth}\begin{align*}
\m_A &= \langle m_1, m_2\rangle, \\
\m_A^2 &= \langle m_1^2, m_1m_2, m_2^2\rangle,\\
\m_A^3 &= \langle m_1^2m_2,m_1m_2^2,m_2^3\rangle,\\
\m_A^4 &= \langle m_1^2m_2^2,m_1m_2^3,m_2^4\rangle,\\
\end{align*}\end{minipage}
\begin{minipage}[t]{0.4\textwidth}\begin{align*}
  \m_A^5 &= \langle m_1^2m_2^3,m_2^5\rangle,\\
\m_A^6 &= \langle m_2^6\rangle,\\
\m_A^7 &= \langle m_2^7\rangle\\
\m_A^8&=0.\end{align*}
\end{minipage}\end{center}
So the first power of $\m_A$ which is zero is $\m_A^8$.
We can see that $m_1\in \Ann_A(\m_A^6)$, since $m_1m_2^6 = m_1^4m_2^2 = 0$, showing that $\Soc^6(A) = \Ann_A(\m_A^6)\neq \m_A^2 = \Rad^2(A)$.
Thus $A$ is not rigid.
\end{example}

\begin{remark}While  \cref{ex:finally} proves that tensor products of field extensions are not generally rigid, it can often happen ``by accident'', even when $E$ does not contain $K$. For example, if $K$ and $E$ are linearly disjoint over $k$, then $K\otimes_k E$ is a field; this happens for example if $E$ is separable and $K$ is purely inseparable. 
  
Or, suppose $K=k(f)$ is a purely inseparable simple extension. Then if we let $r$ be minimal such that $f^r\in E$, the maximal ideal $\m_A$ is the principal ideal generated by the element $x = 1\otimes f^r-f^r\otimes 1$, and rigidity follows easily.
 \end{remark}

The following example helps to motivate our chosen definition of rigidity and shows how base change of a field across a (non-normal) field extension can give rise to an algebra whose regular module has indecomposable modules of different Loewy lengths.
\begin{example}\label{rem:notbottom}
Suppose $p=3$ and let $k = \F_3(t,u)$ be the field of rational functions in $t$ and $u$ over the finite field $\F_3$. 
Let $F = k(t^{1/6} + u^{1/3})$.
Then $F$ has degree $6$ over $k$, and is made up of a Galois extension $E/k$ of degree $2$, where $E = k(t^{1/2})$, together with a purely inseparable extension $F/E$ of degree $3$.
The field $E$ has another purely inseparable extension of degree $3$ which is abstractly isomorphic to $F$, namely $F^* = k(-t^{1/6}+u^{1/3})$.
Let $L = k(t^{1/6},u^{1/3})$ be the compositum of $F$ and $F^*$, so $L/k$ has degree $18$.
The nontrivial Galois automorphism $\gamma: t^{1/2}\mapsto -t^{1/2}$ of $E$ extends in an obvious way to an automorphism of $L$ which swaps the subfields $F$ and $F^*$; we denote this automorphism by $\gamma$ as well.

We claim that the regular module for $A:=F\otimes_k F$ has indecomposable summands of different Loewy lengths.
There are two maximal ideals $\m_1$ and $\m_2$ in $A$: we can realise $\m_1$ as the kernel of the map $x\otimes y\mapsto xy$ with image $F$, and $\m_2$ as the kernel of the map $x\otimes y\mapsto x\gamma(y)$ with image $L$.
Then 
$$
A\cong (F\otimes_E F) \oplus (F\otimes_E F^*).
$$
Both summands are rigid, however the first one has Loewy length $3$ and the second is isomorphic to the field $L$, so has Loewy length $1$.

Some further analysis in this example shows that the minimal field of definition of $\Jac(F_\bark)$ is the field $k':=k(t^{1/3},u^{1/3})$, the subfield of $\gamma$-fixed points in $L$. Hence $F$ is absolutely rigid over $k'$. However, one can also show that $F$ is already absolutely rigid over $k$.
\end{example}

 \subsubsection{Non-connected \texorpdfstring{$G$}{G}}\label{sec:nonconnected}

Having described the field $k_V$ in \cref{thm:mainone} for connected $G$ 
we can finish the proof of that theorem for arbitrary smooth groups.
To begin, we have the following generalisation of \cite[I.6.16]{Jan03} in the smooth case:

\begin{lemma}\label{clifford}
Suppose $G$ is a smooth affine algebraic $k$-group (not necessarily connected) and $N$ is a closed normal subgroup of $G$.
If $V$ is a semisimple $G$-module, then $V$ is a semisimple $N$-module upon restriction.
\end{lemma}

\begin{proof}
The $G_{k_s}$-module $V_{k_s}$ is semisimple by \cref{lem:sepext}(i).
By smoothness the $k_s$-points of $G_{k_s}$ are dense, so we deduce from \cite[I.6.16]{Jan03}
that the $N_{k_s}$-socle of $V_{k_s}$ is a $G_{k_s}$-submodule of $V_{k_s}$ which contains the $G_{k_s}$-socle
of $V_{k_s}$; that is, the restriction of $V_{k_s}$ to $N_{k_s}$ is semisimple.
But then $V$ is semisimple for $N$, by \cref{lem:sepext}(i) again.
\end{proof}

Let $N = G^0$ be the identity component of $G$.
Given an extension $E/k$, we may identify the containment of separable closures $k_s\subseteq E_s$ inside $\overline{E}$.
The restriction of any semisimple $G_{E_s}$-module to $N_{E_s}$ is semisimple, by \cref{clifford}, 
so we can deduce inclusions 
\begin{equation}\label{eq:nest}
\Rad^{l-i}_{N_{E_s}}(V_{E_s}) \subseteq \Rad^{l-i}_{G_{E_s}}(V_{E_s}) \subseteq \Soc^i_{G_{E_s}}(V_{E_s}) \subseteq \Soc^i_{N_{E_s}}(V_{E_s}),
\end{equation}
where $l$ is the Loewy length of $V_{E_s}$ as a $G_{E_s}$-module.
We have:
\begin{lemma}\label{lem:enoughalready}
If $\Rad^{l-i}_{N_{E_s}}(V_{E_s}) = \Soc^i_{N_{E_s}}(V_{E_s})$ for each $i$, then $V_E$ is rigid as a $G_E$-module.
\end{lemma}

\begin{proof}
Under the given hypothesis, all the displayed inclusions become equalities, and these equalities descend to $V_E$, by \cref{lem:sepext}(ii).
This implies that $V_E$ is rigid, by \cref{sumrigid}(iii).
\end{proof}  

To prove \cref{thm:mainone} for disconnected $G$, it remains to establish that the hypothesis of the previous lemma holds with the role of $E$ played by a suitable field $k_V$.
Applying \cref{clifford} to $V$, we see that the restriction of $V$ to $N$ is semisimple.
Let $U$ be one of the simple $N$-summands of the restriction of $V$ to $N$, and base change to $k_s$ so that the $k_s$-points $G_{k_s}(k_s)$ in $G_{k_s}$ are dense.
Let $W := U_{k_s}$, a semisimple $N_{k_s}$-submodule of $V_{k_s}$.
Twisting by an element $g\in G_{k_s}(k_s)$ \cite[I.2.15, I.6.15]{Jan03}, we get another $N_{k_s}$-submodule ${^g}W$.
Summing all such twists we obtain a $G_{k_s}(k_s)$-submodule $X$ of $V_{k_s}$, which must be a $G_{k_s}$-submodule, since $G_{k_s}(k_s)$ is dense in $G_{k_s}$.
But $X$ is also Galois-stable, and so descends to a non-zero $G$-submodule of the simple $V$;
that is, we have $X = V_{k_s}$.
This shows that, as an $N_{k_s}$-module, we have
$$
V_{k_s} = W_1\oplus \cdots \oplus W_r,
$$
where each $W_i = {^{g_i}W}$ for some $g_i\in G_{k_s}(k_s)$.
This remains true after base change to any $E_s$ as above: identifying $G_{k_s}(k_s)\subseteq G_{E_s}(E_s)$, we have $(W_i)_{E_s} = {^{g_i}}(W_{E_s})$, 
and it is clear that the summands have a common Loewy length.
Hence we can apply \cref{sumrigid} to conclude that the terms in the socle and radical series for $V_{E_s}$ (as an $N_{E_s}$-module)
are the sums of the corresponding terms in the summands, and
so the hypothesis of \Cref{lem:enoughalready} holds whenever $W_{E_s} = U_{E_s}$ is rigid. 

Lastly, there is a purely inseparable extension $k_U/k$ attached to the $N$-module $U$ such that
$U_{k_U}$ is absolutely rigid.
If $E/k_U$ is any extension, we have $E_s$ contains $k_U$ also, so $U_{E_s}$ is rigid for $N_{E_s}$, so
\cref{lem:enoughalready} applies by the previous paragraph.
We have now established \cref{thm:mainone} for $G$, 
with $k_V:=k_U$.
\subsubsection{Finite \texorpdfstring{$G$}{G}}
We finish Part I by observing that the analogue of our main result holds for finite algebraic groups.
When $G$ is a finite algebraic group over $k$---i.e., when $k[G]$ is finite rank over the field $k$---we can capture the representation theory of $G$ using the \emph{algebra of measures} $M(G)\cong k[G]^*$ as a replacement for the algebra of distributions $\Dist(G)$,
as explained in \cite[I.8]{Jan03}.
With this replacement made in \cref{sec:morita}, the analogous claim from \cref{eqn:disthoms} holds and the remaining proofs go through \emph{mutatis mutandis} to establish \cref{thm:mainone} for $G$.
This gives a first step towards answering the following general question:

\begin{question}Can one remove from \cref{thm:mainone} the stipulation that $G$ is smooth?\end{question}


%%%%%%%%% PART II

\section{Application of high weight theory}\label{partii}
Let $k$ be a field and $G$ a smooth connected affine $k$-group. The $k$-unipotent radical $\RR_{u,k}(G)$ is the maximal smooth connected normal unipotent $k$-subgroup of $G$, and $G$ is \emph{pseudo-reductive} if
$\RR_{u,k}(G) = 1$. In this part we look to deploy \cite{BS22} and \cite{CPClass} to give a description of $k_V$ in \cref{thm:mainone}. 

\subsection{Preliminaries on the representation theory of pseudo-reductive groups}
We recall that if $k$ is a field and $\U$ is a unipotent $k$-group then we have \cite[Exp.XVII, Prop.~3.2]{SGA3}:
\begin{prop}\label{unipsimp}Let $\U$ be any unipotent $k$-group. Then the only simple $\U$-module is the $1$-dimensional trivial module, $k$.\end{prop}

This implies that if $k$ is a field and $V$ is a simple $G$-module then any normal unipotent subgroup of $G$ acts trivially on $V$. In particular, when $G$ is smooth and connected, the $k$-unipotent radical $\RR_{u,k}(G)$ of $G$ acts trivially on $V$ and in studying $V$ it does no harm to replace $G$ with its maximal \emph{pseudo-reductive} quotient $G/\RR_{u,k}(G)$ and so:

\begin{center}\framebox{In what follows $G$ will always denote a pseudo-reductive algebraic $k$-group over a field $k$.}\end{center}

The next sections recap the main results of \cite{BS22} which describe simple modules for pseudo-reductive groups by means of a high weight theory. 

\subsubsection{Induction} 
In \cite{BS22} simple modules for algebraic $k$-groups are constructed by induction.
We refer the reader to \cite[I.3.3]{Jan03} for the definition, and here just record the key property of \emph{Frobenius reciprocity}  \cite[I.3.4(b)]{Jan03}
for later use. 
Suppose $H$ is a subgroup of $G$. Then for a $G$-module $V$ and an $H$-module $U$
\begin{equation}\label{eq:frob}
\Hom_G(V,\Ind_H^G(U)) \cong \Hom_H(\Res^G_H(V),U),
\end{equation}
where $\Ind_H^G(U)$ is the induced module and $\Res^G_H(V)$ is the $H$-module obtained from $V$ by restriction.

\subsubsection{Pseudo-reductive groups and Levi subgroups}\label{sec:pseudo-prelims} 
 There is a smallest extension $k\subseteq k'\subseteq \bark$ for which $\RR_{u,k'}(G_{k'})_{\bark}=\RR_u(G_{\bark})$, and $k'$ is called the \emph{minimal field of definition} of $\RR_u(G_{\bark})$, where $\RR_u(G_{\bar{k}})$ is the (geometric) unipotent radical of $G$; see \cite[Def. 1.1.6]{CGP15}. By \cite[Prop.~1.1.9]{CGP15} we have that 
the extension $ k'/k$ is finite and purely inseparable.

A pseudo-reductive group $G$ is called \emph{pseudo-split} if it contains a split maximal torus $T$;
in this case, there is a \emph{Levi subgroup} $M$ of $G$ containing $T$.
That, is, there is a split reductive subgroup $M$ of $G$ containing $T$ and such that $G_{\bar{k}} = M_{\bar{k}} \ltimes \RR_u(G_{\bar{k}})$; see \cite[Thm. 3.4.6]{CGP15} or \cite[Thm. 5.4.4]{CPSurv}.




\subsubsection{Weil restriction}\label{sec:Weil} 
We recall some of the important properties of Weil restriction from \cite[\S A.5]{CGP15}.	 
If $B\to B'$ is a finite flat map of noetherian rings, and $X'$ a quasi-projective $B'$-scheme, one may define the Weil restriction 
$X:=\R_{B'/B}(X')$. Then $X$ is a $B$-scheme of finite type satisfying the universal property
\[X(A)=X'(B'\otimes_B A),\]
for $A$ any $B$-algebra. If $B\to C$ is a further map of rings, then \cite[A.2.7]{Oes84} gives 
\begin{equation}\label{eq:weilandback}X_C\cong \R_{(B'\otimes_B C)/C}(X').\end{equation}

A key fact is that Weil restriction is right adjoint to base change along $\Spec(B)\to \Spec(B')$. That is to say that there is a bijection
\begin{equation}\label{weiladj}\Hom_B(Y,\R_{B'/B}(X'))\cong \Hom_{B'}(Y_{B'},X'),\end{equation}
which is natural in $X'$ and the $B$-scheme $Y$. 
One situation is particularly important below. 
If $X'=Z_{B'}$ for a $B$-scheme $Z$ then taking $Y=Z$ in \cref{weiladj}, one has the identity map on the right-hand side, giving a canonical map $Z\to \R_{B'/B}(X')$; \cite[A.5.7]{CGP15} implies that this map is a closed immersion provided $\Spec(B')\to \Spec(B)$ is surjective (which is true if $B$ is a field and $B'$ is non-zero, since then $\Spec(B)$ is a single point). 

In case $X'=G'$ is a $B'$-group, we find $G:=X$ is a $B$-group. 
When $B=k$ is a field, and $B'=k'$ is a nonzero finite reduced $k$-algebra, then $G'$ is pseudo-reductive whenever $G$ is (connected) reductive.  
If $G'$ is defined over $k$ and we choose a $k$-descent $H$ of $G'$, then the remarks above show that $H$ embeds as a canonical subgroup in $G$; in particular this holds in the case that $G' = G_{k'}$ for some $k$-group $G$, giving a canonical embedding of $G$ in $\R_{k'/k}(G_{k'})$ 
which we refer to below as the ``natural copy of $G$ coming from adjunction''.

We recall also a feature of Weil restriction across a separable extensions $l/k$. Let $X$ be an affine scheme of finite type over $l$ and $\Gamma$ the Galois group of $\hat l/k$, where $\hat l$ is a normal closure of $l$ in $\bar l$. Denote by $\Gamma_l$ the subgroup of $\Gamma$ fixing $l$. Then \begin{equation}\label{eqn:weilsep}\R_{l/k}(X)_{\hat l}\cong \prod_{\gamma}(\gamma(X_{\hat l})),\end{equation} where $\gamma$ runs over a transversal of $\Gamma_l$ in $\Gamma$.




\subsubsection{Simple modules for pseudo-split groups}\label{sec:BS}
We summarise the main results from \cite{BS22}.
Suppose $G$ is a pseudo-split pseudo-reductive $k$-group, with maximal torus $T$ and Cartan subgroup $C = Z_G(T)$. 
There is a Levi $k$-subgroup $M$ of $G$ containing $T$, and having chosen a Borel subgroup of $M$ (which, for technical reasons should correspond to the negative roots), we get a system of positive roots, from which we can define a set of dominant weights $X(T)_+$ for $T$.
The following is a portmanteau theorem from \cite[Thm.~1.2, Thm.~3.1]{BS22}.

\begin{theorem}\label{thm:fromBS}
Suppose $\lambda\in X(T)_+$. Let $Q_G(\lambda):= \Ind_M^G(L_M(\lambda))$.
\begin{itemize}
\item[(i)] The socle of $Q_G(\lambda)$ is a simple module, denoted $L_G(\lambda)$.
\item[(ii)] The assignment $\lambda \to L_G(\lambda)$ gives a one-one correspondence between the dominant weights $X(T)_+$ and the simple $G$-modules.
\item[(iii)] The highest weight space $L_G(\lambda)_\lambda$ of $L_G(\lambda)$ is a $C$-module isomorphic to $L_C(\lambda)$.
\item[(iv)] The restriction $\Res^G_M(L_G(\lambda))$ is isotypic and semisimple, and hence 
$$
\dim(L_G(\lambda)) = \dim L_C(\lambda)\cdot \dim L_M(\lambda).\label{thm:fromBS:dim}
$$
\end{itemize}
\end{theorem}

\begin{example}\label{rem:notabs}
It is noted in \cite[Rem 4.7(i)]{BS22} that the modules $L_G(\lambda)$ are rarely absolutely semisimple.
The most basic example is as follows: let $k$ be an imperfect field of characteristic $p$ and $G = \R_{k'/k}(\Gm)$ be the Weil restriction of the multiplicative group
across a purely inseparable extension $k'/k$ of degree $p$. (Then $G$ is pseudo-split and pseudo-reductive---see \cite[1.1.3]{CGP15}.) 
For $(r,p)=1$, the simple module $L_G(r)$ can be realised as the Weil restriction $\R_{k'/k}(V)$, where $V\cong k'$ is a $1$-dimensional vector space on which $\Gm$ acts with weight $r$; in other words, if $g\in \Gm(A)=A^\times$, then $g\cdot v\mapsto g^rv$ for any $v\in V(A)\cong A$.  Noting that $G(k)=\Gm(k')=(k')^\times$, one sees that $G(k)$ has one orbit on the non-zero elements of $L_G(r)$ and so $L_G(r)$ is $p$-dimensional and simple, but $L_G(r)_{k'}$ 
is reducible and indecomposable as a module for $G_{k'}$. 
(To see this latter statement, if we write $k' = k(a)$ with $k$-basis $1,a,\ldots,a^{p-1}$, then one can realise the elements of $G$ in their action on $L_G(r)$ explicitly as 
$p\times p$ matrices. Over $k'$ these matrices are trigonalisable.)

When $r=1$, this describes (the Weil restriction of) the natural action of $\Gm$ on $\Ga$ coming from scalar multiplication in $k'$.
See \cref{cor:sepclosedsimple} below for a contrasting result for separable base changes of simple modules.
\end{example}


\subsubsection{Simple modules and the map \texorpdfstring{$i_G$}{iG}}\label{sec:i_G}
Keep notation from the previous section, so $G$ is a pseudo-split pseudo-reductive group with maximal torus $T$, Cartan $C = Z_G(T)$, and Levi subgroup $M$ containing $T$.
Let $K$ be the minimal field of definition of the unipotent radical of $G$, set $U = \RR_{u,K}(G_{K})$, and let $\G:= \R_{K/k}(M_{K})$.
The quotient $\pi:G_{K} \to M_{K}$ induces a map $i_G:G \to \G$ as follows, see also \cite[Eq. (1.3.1)]{CGP15}.
We can write $G_{K} = M_{K}\ltimes \RR_{u,K}(G_{K})$, and therefore we have 
$$
\R_{K/k}(G_{K}) = \G\ltimes \R_{K/k}(U),
$$
and this group contains a natural copy of $G$ coming from adjunction.
The Weil restriction $\R_{K/k}(\pi)$ is just the quotient map of the displayed semidirect product by the smooth normal 
unipotent $k$-subgroup $\R_{K/k}(U)$, and thus we can define a map 
\begin{equation}\label{eq:i_G}
i_G:G \to \G
\end{equation} 
as the composition of the embedding of $G$ in the semidirect product with the quotient $\R_{K/k}(\pi)$.
This map plays a crucial role in the structure theory of pseudo-reductive groups developed in \cite{CGP15}.
We note that $i_G$ often has trivial kernel, but not always: see \cite[Ex. 1.6.3, Ex. 5.3.7]{CGP15} for examples of this.

It is proved in \cite[Prop 7.1.3(ii)]{CPSurv} that $\ker i_G$ is unipotent (with no non-trivial smooth and connected subgroups) and the smooth connected image $i_G(G)$ is pseudo-reductive. When the intersection of $\ker i_G$ with a Cartan subgroup of $G$ is trivial, we say $G$ is of \emph{minimal type}.


It is also shown in \cite[Thm. 1.6.2(2)]{CGP15} that when $H = \R_{K/k}(H')$ is the Weil restriction of a reductive $K$-group $H'$ across the purely inseparable extension $K/k$, then the map $i_H$ is an isomorphism.
From this, it is not hard to deduce that  the map $i_{i_G(G)}$ is nothing other than the inclusion $i_G(G)\hookrightarrow \G$. 

Together with \cref{unipsimp}, this implies:
\begin{lemma}\label{lem:i_G}
  The action of $G$ on $L_G(\lambda)$ factors through $G\to i_G(G)$. Thus we may assume $G=i_G(G)\subseteq \G$; in particular $G$ is of minimal type.
\end{lemma}

We make use of this observation in \cref{prop:splitend} below, comparing the simple modules of $G$ and $\G$.
%In particular, when we consider below the restriction of a $\G$-module to $G$, or the induction from $G$ to $\G$, we should really be considering the image $i_G(G)$ of $G$ inside $\G$, but with the previous point made this should not cause any confusion.

\subsection{Endomorphisms of simple modules: pseudo-split case}\label{sec:endsimple1} 
In this and the following sections, we show how to pin down the endomorphism algebra $D = \End_G(V)$ more precisely, which allows us to be quite explicit about the field $k_V$ in \cref{thm:mainone}. In this section $G$ is pseudo-split with  a split maximal torus $T$ that is contained in a Levi subgroup $M$ and $C=Z_G(T)$ is a Cartan subgroup. We let $K$ denote the minimal field of definition of the unipotent radical of $G$ thence get the canonical map $i_G:G\to\G:=\R_{K/k}(M_K)$, which restricts to $i_C:C\to\C:=\R_{K/k}(T)$. Evidently $T$ is a Levi subgroup of $C$ and $\C$; as $M$ is of $G$ and $\G$. 

It is helpful to relate the simple modules of $G$ and $\G$, building on work in \cite{BS22}.  Recall the following from \cite[Sec.~1]{BS22}:
\begin{defn}\label{def:Klambda}
Given any $T$-weight $\lambda = (\lambda_1,\ldots,\lambda_r)\in\Z^r$, for each $i$ we can form the subfield $K(\lambda_i)$ of $K$ 
generated by $k$ together with $(K)^{\lambda_i}$. 
As $K/k$ is purely inseparable, taking $\lambda_i = p^e\mu$ with $\mu$ coprime to $p$, we have $K(\lambda_i) = K(p^e)$.
Let $K(\lambda)$ denote the purely inseparable subfield of $K$ generated by $k$ and the $K(\lambda_i)$.
Given the data of $K/k$ and the weight $\lambda$, we call $K(\lambda)$ the \emph{field attached to $\lambda$}.
\end{defn}

As explained in \cite[Thm.~5.8]{BS22}, the simple module $L_{\C}(\lambda)$ identifies with the field $K(\lambda)$; in such a way that the action of $\C$ on $L_\C(\lambda)$ factors through the natural action of $\R_{K(\lambda)/k}(\Gm)$ on $\R_{K(\lambda)/k}(\Ga)$ via a surjection $\C\to \R_{K(\lambda)/k}(\Gm)$. 


\begin{prop}\label{prop:splitend}
Suppose $\ker i_G=1$ and identify $G=i_G(G)$ as a subgroup of $\G$. For any $T$-weight $\lambda$ we have:
\begin{enumerate}
\item The simple module $L_C(\lambda)$ can be identified with an intermediate extension $k\subseteq K_C(\lambda) \subseteq K(\lambda)$, 
and the action of $C$ on $L_C(\lambda)$ factors through the natural action of $\R_{K_C(\lambda)/k}(\Gm)$ on $\R_{K_C(\lambda)/k}(\Ga)$ by scalar multiples.
\item More generally,  we have an isomorphism of $G$-modules $L_G(\lambda)\cong \R_{K_C(\lambda)/k}(L_M(\lambda)_{K_C(\lambda)})$, with $G$
acting through $\R_{K_C(\lambda)/k}(\GL(L_M(\lambda)_{K_C(\lambda)}))$.
Further, by restriction, we have
$$
\Res^{\G}_G(L_{\G}(\lambda)) \cong L_G(\lambda)\oplus\cdots\oplus L_G(\lambda),
$$
with $[K(\lambda):K_C(\lambda)]$ summands on the right-hand side.\vspace{5pt}
\item \hspace{0.3\textwidth}$\End_G(L_G(\lambda))\cong K_C(\lambda).$\label{prop:splitend:end}
\end{enumerate}
\end{prop}

\begin{proof}
By \cite[Lem.~5.10]{BS22},  any subgroup of $\C$ stabilising a proper $k$-subspace of $K(\lambda)$ 
actually stabilises a proper subfield of $K(\lambda)$---in fact, the proof shows that $1\in K(\lambda)$ generates a canonical $C$-submodule under $C$ that we denote $K_{C}(\lambda)$; furthermore, \emph{loc.~cit.} shows that $C$ acts on it through the Weil restriction $\R_{K_C(\lambda)/k}(\Gm)$.
This proves (i).

When $G$ is commutative, so $G=C$ and $M=T$, then (ii) is an easy application of (i) with \cref{clifford}; then (iii) follows since the $C$-action commutes with the the multiplicative structure of $K_C(\lambda)$ and any endomorphism $\phi:K_C(\lambda)\to K_C(\lambda)$ is completely determined by the image of $1\in K_C(\lambda)$.

With this in hand, we may prove  (iii) for general $G$. Note that any $G$-module homomorphism $\phi:L_G(\lambda)\to L_G(\lambda)$ must stabilise the high weight space $L_G(\lambda)_\lambda$ and is completely determined by what happens to it, since its vectors generate the whole of $L_G(\lambda)$---even under a Levi subgroup $M\subseteq G$.
Hence $\phi$ is determined by its restriction to $L_G(\lambda)_\lambda\cong L_C(\lambda)$, which by the above shows that we have an inclusion $\End_G(L_G(\lambda))\subseteq K_C(\lambda)$.

On the other hand, any $G$-homomorphism $L_G(\lambda) \to Q_G(\lambda)$ must land in the simple socle $L_G(\lambda)$, so we have $\End_G(L_G(\lambda))\cong \Hom_G(L_G(\lambda),Q_G(\lambda))$.
But now we can use Frobenius reciprocity \eqref{eq:frob}:
$$
\Hom_G(L_G(\lambda),Q_G(\lambda)) \cong \Hom_M(\Res^G_M(L_G(\lambda)),L_M(\lambda)).
$$
Since the restriction $\Res^G_M(L_G(\lambda))$ is isomorphic to a direct sum of copies of $L_M(\lambda)$, we can conclude that $\End_G(L_G(\lambda)) \cong k^t$, where $t$ is the number of summands. 
But the dimension formula in \cref{thm:fromBS}(iv) says that 
$t = \dim_k L_C(\lambda) = \dim_k K_C(\lambda)$, so we are done.

Finally to complete the proof for (ii), let $M\subseteq G\subseteq \G$ be a Levi subgroup. Then we have
\begin{align*}\Res^\G_M(L_\G(\lambda))&\cong L_M(\lambda)^{\dim L_\C(\lambda)}&{\text{(Thm.~\ref{thm:fromBS}(iv))}}\\
&\cong L_M(\lambda)^{[K(\lambda):k]}.
\end{align*}

Since also $\Res^G_M(L_G(\mu))\cong L_M(\mu)^{[K_C(\mu):k]}$ by \cref{thm:fromBS}(iv) applied to $G$, it follows that all the $G$-composition factors of $\Res^\G_G(L_\G(\lambda))$ are isomorphic to $L_G(\lambda)$ and there are $$[K(\lambda):k]/[K_C(\lambda):k]=[K(\lambda):K_C(\lambda)]$$ of them. It remains to show that $\Res^\G_G(L_\G(\lambda))$ is semisimple.

For that we calculate the dimension of $\Hom_G(\Res^\G_G(L_\G(\lambda)),L_G(\lambda))$, which is equal to that of $\Hom_G(L_\G(\lambda),\Ind_G^\G(L_G(\lambda)))$ by Frobenius reciprocity \eqref{eq:frob}. Now $$Q_\G(\lambda)=\Ind_M^\G(L_M(\lambda))\cong \Ind_G^\G(\Ind_M^G(L_M(\lambda)))=\Ind_G^\G(Q_G(\lambda))$$ by transitivity of induction. 
As $Q_\G(\lambda)$ has a simple socle, we get
\[\End_\G(L_\G(\lambda))\cong \Hom_\G(L_\G(\lambda),Q_\G(\lambda))\cong \Hom_G(\Res^\G_G(L_\G(\lambda)),L_G(\lambda))\]
by Frobenius reciprocity again. 
Part (iii) of this lemma identifies the left-most term as having the structure of $K(\lambda)$ as a $k$-vector space. 
Meanwhile, the right-most term has $k$-dimension equal to $\dim\End_G(L_G(\lambda))=[K_C(\lambda):k]$ times the multiplicity of $L_G(\lambda)$ in the head of $\Res^\G_G(L_\G(\lambda))$. Therefore that multiplicity must be $[K(\lambda):K_C(\lambda)]$ and we are done.
\end{proof}

\begin{remark}
  Note that $L_G(\lambda)$ is \emph{canonically} embedded in $L_\G(\lambda)$ since $L_C(\lambda)$ is canonically embedded in $L_\C(\lambda)$ due to the inclusion of fields $K_C(\lambda)\subseteq K(\lambda)$. More specifically, we see $L_G(\lambda)$ as the image under $M$ of the subfield $K_C(\lambda)$ in the highest weight space of $L_{\G}(\lambda)$---which we have identified with $K(\lambda)$.
\end{remark}

Evidently $K_C(\lambda)$ is as natural a finite purely inseparable extension of $k$  attached to $V=L_G(\lambda)$ as one might reasonably ask for.
By \cref{prop:splitend}, $D:=\End_G(L_G(\lambda))\cong K_C(\lambda)$ is also a finite purely inseparable extension of $k$, and so by \Cref{lem:mfod}(i) it is 
itself the minimal field of definition of $D$.
The following is immediate.
\begin{cor}\label{cor:splitrigid}
  With terminology as above, $k_V:=K_C(\lambda)$ satisfies the conclusion of \cref{thm:mainone}.
\end{cor}


\begin{example}\label{ex:klambda}
Easy examples show that the fields $K_C(\lambda)$ and $K(\lambda)$ in \cref{prop:splitend} can be different.
Let $k = F(a,b)$ be the field of rational functions in two indeterminates over a field $F$ of characteristic $p$,
and let $K = k(a^{1/p},b^{1/p})$, a purely inseparable extension of $k$ of degree $p^2$.
Set $s = a^{1/p}$ and $t=b^{1/p}$.
Let $C = \R_{k(s)/k}(\Gm)\times \R_{k(t)/k}(\Gm)$, a commutative pseudo-split pseudo-reductive group with maximal split torus $T=\Gm\times\Gm$.
Then $K$ is the minimal field of definition of $\RR_u(C_\bark)$, so the group $\C = \R_{K/k}(T_K) \cong \R_{K/k}(\Gm)\times \R_{K/k}(\Gm)$,
with the factors of $C$ sitting naturally inside the factors of $\C$.

Consider the three $T$-weights $(1,0)$, $(0,1)$ and $(1,1)$.
The corresponding modules for $\C$ are all isomorphic as $k$-vector spaces to the field $K$ itself: since there are no nontrivial powers of $p$ appearing, we have $K(\lambda) = K$ for each of the three choices of $\lambda$.
On the other hand, we have $K_C(1,0) = k(s)$, $K_C(0,1) = k(t)$ and $K_C(1,1) = k(s,t) = K$,
and so we see that the field $K_C(\lambda)$ does depend on the weight and on the group $C$.
\end{example}




\begin{remark}\label{rem:finally}
Recall \cref{ex:finally}, which gives an example of a field $k$ and purely inseparable extensions $K/k$ and $E/k$ such that $K\otimes_k E$ is not rigid.
By setting $G = \R_{K/k}(\Gm)$ and $V = L_G(1)$, we obtain an example of a pseudo-split pseudo-reductive group $G$, a simple $G$-module $V$ with $\End_G(V) = K$, 
and an extension $E/k$ such that $V_E$ is not rigid.
This shows that we cannot hope that simple $G$-modules are absolutely rigid, even when $G$ is pseudo-split. 
Similar examples can be constructed replacing $\Gm$ with other split reductive groups, 
and will occur whenever we have purely inseparable extensions whose tensor product is not rigid.
\end{remark}

It is a fact---\!\!\cite[II.2.9]{Jan03}---that simple $G$-modules for split reductive $G$ are all defined over the relevant prime fields. With knowledge of the endomorphism ring in hand, we can give the generalisation to pseudo-split pseudo-reductive groups. 
At the same time we observe that the simple $G$-modules are absolutely indecomposable.

\begin{cor}\label{cor:sepclosedsimple}
Let $G$ be pseudo-split and $\lambda\in X(T)_+$ with $K_C(\lambda)$ the field associated with $L_G(\lambda)_\lambda\cong L_C(\lambda)$. 
\begin{itemize}
\item[(i)] We have $L_G(\lambda)_E = L_{G_{E}}(\lambda)$ for any field extension $E/k$ which is linearly disjoint from $K_C(\lambda)$---for example if $E/k$ is separable.
\item[(ii)] We have $L_G(\lambda)_E$ is isotypic with simple socle and head---hence indecomposable---for every field extension $E/k$.
\end{itemize}
\end{cor}

\begin{proof} 
(i). By \cref{prop:splitend} we have \[\End_{G_{E}}(L_G(\lambda)_{E})\cong \End_G(L_G(\lambda))_{E}\cong K_C(\lambda)\otimes_k E,\]
and if $K_C(\lambda)$ and $E$ are linearly disjoint, then $K_C(\lambda) \otimes_k E$ is a field. 
It follows that $L_G(\lambda)_{E}$ is simple.

(ii). By the results of \Cref{sec:morita}, the submodule structure of $L_G(\lambda)_E$ is controlled by the ideal structure of $K_C(\lambda)\otimes_k E$.
Since $K_C(\lambda)/k$ is purely inseparable, this tensor product is a Gorenstein local ring which implies the statement.
\end{proof}



\subsection{Endomorphisms of simple modules: general case}\label{sec:endsimple2}
We drop the assumption that $G$ is pseudo-split. Then for $V$ a simple $G$-module, $D:=\End_G(V)$ will no longer be a field in general, let alone a purely inseparable extension of $k$. We  elucidate the structure of $D$ taking inspiration from \cite{Tit71}.
A classification of the possible isomorphism classes of $D$ that could occur would subsume many difficult open questions about the Brauer groups of fields and we do not try to tackle this here. Instead, we assume we know the action of the absolute Galois group on (the Dynkin diagram of) $G$ and $V$. Then we are able to describe $D$ through its base change to a suitable separable extension---after which $G$ becomes pseudo-split and we can deploy the highest weight theory discussed in the last section. 
In particular we can calculate the dimension of $D$ based on this data.

Let $T$ be a maximal torus of $G$ and let $E/k$ be a finite Galois extension such that $S:=T_E$ is split, and hence $G_E$ is pseudo-split.
By choosing a Borel subgroup of a Levi subgroup of $G_E$ containing $S$, we can fix a system of dominant weights $X(S)_+$ in the weight lattice for $S$.
Let $\Gamma = {\rm Gal}(E/k)$, and let $W$ be the Weyl group of $G_E$.

We recall two actions of $\Gamma$ on the weights of $S$, see also \cite[Sec. 3.1]{Tit71}. 
The first arises from base change: given $\gamma\in \Gamma$ we can form the base change along $\gamma$ of the torus $S$,
giving a torus ${^\gamma}S$ defined functorially
by the formula ${^\gamma}S(A) = S(A\otimes_\gamma E)$ for each $E$-algebra $A$.
Since $k$ is fixed by $\gamma$, the tori $S$ and ${^\gamma}S$ have the common $k$-descent $T$, and so they are naturally isomorphic---we can identify ${^\gamma}S$ with $S$.
Under this identification, when we base change a character $\lambda$ of $S$ along $\gamma$ we obtain a new character ${^\gamma}\lambda$ of $S$.
Typically this will not preserve the dominant weights, but note that for each $\gamma\in \Gamma$ there is a unique $w\in W$ such that $w({^\gamma}X(S)_+)\subseteq X(S)_+$, and for $\lambda\in X(S)$ we set
\begin{equation}\label{eq:dot}
\gamma\cdot\lambda:=w({^\gamma}\lambda).
\end{equation}
Note that this action respects the partial order on weights, since the system of positive roots corresponding to the choice of Borel subgroup above must also be preserved.

Let $V$ be a simple $G$-module.
Then $\Gamma$ acts semilinearly on $V_E = V\otimes_k E$ via its action on $E$;
denote this action by $v\mapsto \gamma(v)$. 
An $E$-subspace of $V_E$ has a $k$-form if and only if it is $\Gamma$-stable.
Note also that if $v\in V_E$ is a vector of $T_E$-weight $\lambda$,
then $\gamma(v)$ has weight ${^\gamma}\lambda$.
Since $V$ is simple as a $G$-module, $V_E$ is semisimple as a $G_E$-module by \cref{lem:sepext}, and so $V_E$ is
the sum of certain simple $G_E$-modules.
Let $\Lambda = \{\lambda=\lambda_1,\ldots,\lambda_r\}\subset X(S)_+$ be the set of highest weights occurring.
The following proof is based on that in \cite[Sec. 7.6]{Tit71}.

\begin{lemma}\label{lem:Vlambda}
Keep the notation above.
Further, for each $\lambda\in \Lambda$, let $V\{\lambda\}$ denote the sum of the simple submodules of $V_E$ isomorphic to $L_{G_E}(\lambda)$. Then:
\begin{itemize}
\item[(i)] $\Lambda$ forms a single $\Gamma$-orbit in $X(S)_+$;
\item[(ii)] $V_E$ is the direct sum of the $V\{\lambda\}$;
\item[(iii)] for each $\lambda_i\in\Lambda$, the multiplicity of $L_{G_E}(\lambda_i)$ in $V_{E}$ is a fixed integer, $d$.
\end{itemize}
\end{lemma}

\begin{proof}
Let $U$ be a simple $G_E$-submodule of $V$ isomorphic to $L_{G_E}(\lambda)$ and let $\gamma\in \Gamma$. The subspace $\gamma(U)$ is still a simple submodule, and any weight has the form ${^\gamma}\mu$ for a weight $\mu$ occurring in $U$.
Since the set of weights of $\gamma(U)$ is stable under the action of the Weyl group $W$, in fact any weight of $\gamma(U)$ has the form $\gamma\cdot\mu$ for
$\mu$ occurring in $U$. 
We have observed above that the ordering of weights is preserved by the map $\mu\mapsto \gamma\cdot\mu$,
so we can conclude that the module $\gamma(U)$ is isomorphic to $L_{G_E}(\gamma\cdot\lambda)$.

Now let $X$ be the (non-trivial) submodule of $V_E$ generated by the $\gamma(U)$ for $\gamma\in \Gamma$.
Then $X$ is $\Gamma$-stable, and hence has a $k$-form, which corresponds to a non-trivial $G$-submodule of $V$.
Since $V$ is simple, we conclude that $X = V_E$.
This proves (i), and (ii) follows because $V_E$ is semisimple.
For (iii), the above considerations imply that $\gamma(V\{\lambda\}) = V\{\gamma\cdot\lambda\}$ for each $\lambda\in \Lambda$ and $\gamma\in \Gamma$.
Hence the multiplicity of $L_{G_E}(\lambda)$ as a summand of $V\{\lambda\}$ must equal the multiplicity of $L_{G_E}(\gamma\cdot\lambda)$ as a summand of $V\{\gamma\cdot\lambda\}$.
\end{proof}

We can now describe the minimal field of definition of $\Jac(D_\bark)$ which serves the role of the field $k_V$ in \Cref{thm:mainone}.
According to the previous lemma and \Cref{cor:sepclosedsimple}(i), when we further extend to $k_s$ we get a decomposition 
\begin{equation}\label{ksdecomp}
V_{k_s} = \bigoplus_{\lambda\in \Lambda} L_{G_{k_s}}(\lambda)^d.
\end{equation}
For each $\lambda\in \Lambda$, denote by $K_{C,s}(\lambda)$ the purely inseparable extension of $k_s$ that identifies with the high weight space of $L_{G_{k_s}}(\lambda)$.
The compositum $K$ of the $K_{C,s}(\lambda)$ inside $\bar{k}$ is stable under the absolute Galois group, since the action of the Galois group permutes the summands, and hence permutes their high weight spaces.
This means that $K/k_s$ descends to a purely inseparable extension $k'/k$. 

\begin{theorem}
  With the above notation, $k'$ is the minimal field of definition of $\Rad_{G_\bark}(V_\bark)$ as a module. We may identify $k'$ with the minimal field of definition of $\Jac(D_\bark)$. It follows that we may take $k_V=k'$ in \cref{thm:mainone}.
  \end{theorem}
  \begin{proof}
    Let $\ks$ be the minimal field of definition of $\Jac(D_\bark)$. From the Morita equivalence in \cref{prop:REDE}: for any field extension $E/k$ we have $D_E/\Jac(D_E)$ is absolutely semisimple if and only if $\Soc_{G_E}(V_{E})$ is absolutely semsimple; and since duality preserves the simple modules this happens if and only if $V_E/\Rad_{G_E}(V_E)$ is absolutely semisimple. Therefore $\ks$ identifies with the minimal field of definition of $\Rad_{G_\bark}(V_\bark)$.

Now \eqref{ksdecomp} implies
  \begin{equation}
    D_{k_s} \cong \bigoplus_{i=1}^r M_d(K_i).  \label{eq:ksdecomp}
  \end{equation}
  Since $K_i/k_s$ is purely inseparable, the minimal field of definition of $\Jac(M_d(K_i)_\bark)$ over $k_s$ is $K_i$ itself (\Cref{lem:mfod}(i)), and hence the minimal 
  field of definition of $\Jac(D_\bark)$ over $k_s$ is the compositum of these fields; that is, this minimal field of definition is $K\cong k'\otimes_k k_s$. Hence $k'\otimes_k k_s =\ks\otimes_k k_s$ as subfields of $\bark$ and so $k'=\ks$ as required. 
  
  For the last sentence, we apply \cref{thm:Drigid}.
  \end{proof}

In the context of \cref{rem:noncan}, the next theorem observes how restrictive the demand is for a division ring $D$ to be the endomorphism algebra $\End_G(V)$ for $V$ a simple $G$-module. In other words, the division algebras arising through \cref{lem:schur} come from a particularly special class.

\begin{theorem}\label{thm:torussplit}
  Let $D:=\End_G(V)$ for $V$ a simple $G$-module. Then there is a unique $p$-splitting field $E/k$ for $D$, and $E/k$ is Galois. This field identifies as the splitting field of a maximal torus $T$ in the image $\barG$ of $G$ in $\GL(V)$
\end{theorem}
\begin{proof} The existence of the unique minimal extension $E/k$ and its property of being Galois is a consequence of the discussion in \cite[\S8.12]{Bor91}.
\footnote{See \href{https://mathoverflow.net/questions/142801/splitting-field-of-a-torus}{MathOverflow 142801} for an in-depth analysis of this point.}
  
  Now $\barG_E$ is pseudo-split and pseudo-reductive, so the endomorphism algebra $D_E:=\End_{G_E}(V_E)$ is $p$-split by \cref{lem:Vlambda} and \cref{prop:splitend}\cref{prop:splitend:end}. For the converse, we assume $E$ is such that $D_E$ is $p$-split. 
  
  For some $d$ and $r$ we have $D_E=\Mat_d(E_1)\times\dots\times \Mat_d(E_r)$ where the $E_i$ are purely inseparable over $E$. We have $D_E$ has $r$ simple right modules, which remain simple after separable extension. By Morita equivalence, the same is true of $V_E$. Say $V_E\cong V_1^d\oplus \dots\oplus V_r^d$ with each $V_i$ being $k_s$-simple. Now $V_i$ identifies with some $L_{G_{k_s}}(\lambda)$ over $k_s$, and it descends to $k$, giving a high-weight module $L_G(\lambda)$ on which $T$ acts completely reducibly. Iterating over the composition factors of $V_E$ we see that the image of $T$ in $\barG$ is split.  
\end{proof}


\begin{example}
  Let $k:=\F_p(t)$, $E:=\F_{p^2}(t)$, $k':=\F_p(t^{1/p})$, $F:=\F_{p^2}(t^{1/p})=E\otimes_k k'$. Then $\Gal(E/k)=\langle \gamma\rangle=\Gal(F/k')$, say. Now for $n\geq 2$, denote by $\G$ the (reductive) $k$-group $\R_{E/k}(\SL_n)$ and take $G'$ to be the subgroup scheme of $\G$ given by $G'(A):=\{x\in \G(A)\mid x^\intercal\gamma(x)=1\}$ for any commutative $k$-algebra $A$. 

  The matrices $G'(k)$ describe the non-split reductive $k$-subgroup $\SU_n$;\footnote{In fact $G'$ is \emph{quasi-split} as it contains the descent of a Borel subgroup of $\GG$.} this means that the Galois group attached to $E/k$ induces a non-trivial involution of the Dynkin diagram of ${G'}_E$---see \cite[\S24.f]{Milne17}. Let $G:=\R_{k'/k}({G'}_{k'})$. We have that $G$ is pseudo-reductive and has a canonical $k$-subgroup $G'$ which is evidently a Levi subgroup for $G$. It is well-known, and easy to calculate that the action of $\langle\gamma\rangle=\Gal(E/k)\cong C_2$ on $X(T)_+$ is given by $\gamma: \lambda=(\lambda_1,\dots,\lambda_{n-1})\mapsto \gamma\cdot\lambda=(\lambda_{n-1},\dots,\lambda_1)$. In particular if $\lambda=\varpi_1$ is the first fundamental dominant weight---so that $L_{G_E}(\varpi_1)$ is the natural $pn$-dimensional representation of the pseudo-split $G_{E}$---then there is a simple $G$-module $V=L_G\{\varpi_1,\varpi_{n-1}\}$ of dimension $2pn$ over $k$ such that $V_E\cong L_{G_E}(\varpi_1)\oplus L_{G_E}(\varpi_{n-1})$. So $\End_{G_E}(V_E)\cong F\oplus F$ and is commutative. Since $F$ has an $E/k$-form $k'$, we see that the minimal field of definition of   $\Jac_{\bark}(D_\bark)$ is $k'$. Since the unique non-trivial field extension of $k$ contained in $k'$ is $k'$ itself, this shows that $Z(D)$---hence also $D$ and $V$---are all absolutely rigid.
\end{example}

\subsection{Applying the Conrad--Prasad classification}
We finish by showing how the linear algebra data determining the pseudo-reductive groups of minimal type can be used to identify $K_C(\lambda)$. Recall the headline result of the Conrad--Prasad classification, \cite[Thm.~9.2.1]{CPClass}.

\begin{theorem}[Structure theorem]
  Let $G$ be a pseudo-reductive group over a field $k$. Then $G$ is generalised standard if and and only if it is locally of minimal type.
\end{theorem}

We explain the terminology. If $T$ is a split maximal torus in an affine $k$-group $H$, and $a$ is a root of $T$ (a non-zero weight on $\Lie(H)$), then one may define the root groups $H_a$ of $H$ by using the limits of cocharacters adapted to the root $a$---see \cite[\S2.3]{CGP15}. If $H$ is also pseudo-reductive, and $a$ is non-divisible, then the subgroup $H_{\pm a}$ generated by $H_a$ and $H_{-a}$ is pseudo-simple and pseudo-split with $A_1$ or $BC_1$ root system. Now, $G$ is \emph{locally of minimal type} if for a maximal $k$-torus $T$, the subgroup $(G_{k_s})_{\pm a}$ admits a pseudo-simple central extension of minimal type. If $G$ is of minimal type, then it is must be locally of minimal type---\hspace{1sp}\cite[\S4.3]{CPClass}. Since $\ker(i_G)$ is a unipotent group and we are interested in describing simple $G$-modules, we may assume (c.f.~\cref{lem:i_G}):
\begin{center}\framebox{From now on $\ker(i_G)=1$; in particular, $G$ is locally of minimal type.}\end{center}

Hence the theorem above tells us that $G$ is generalised standard, and so we now explain that construction. First suppose $k'/k$ is a non-zero finite reduced $k$-algebra $k'$ and let $G'$ be a $k'$-group. Then we have $k'\cong k_1'\times\dots\times k_r'$ as a product of factor fields with say $G_i'$ the fibre of $G'$ over $k_i'$. Now $(G',k'/k)$ is said to be a \emph{primitive pair}\footnote{This is \cite[Defn.~9.1.5]{CPClass}---see the comments around \textit{loc.~cit.~}for the definition of each type given.} if each fibre $G'_i$ is one of the following:
\begin{enumerate}\item a connected semisimple, absolutely simple, and simply connected $k_i'$-group;
  \item 
  \begin{enumerate}
  \item a basic exotic group of type $G_2$ ($p=3$) or $F_4$ ($p=2$);
  \item a generalised basic exotic group of type $B$ ($p=2$) [note that this contains the non-standard ];
  \item a generalised basic exotic group of type $C$ ($p=2$);
  \item a rank-$2$ basic exceptional group of type $B_2$ ($p=2$);
  \end{enumerate}
  \item a minimal-type absolutely pseudo-simple with a non-reduced root
  system over the separable closure of $k_i'$ and root field equal to $k_i'$ ($p=2$).\label{nonredcase}
\end{enumerate}
If $C$ denotes a Cartan subgroup of $G$ then the $k$-group functor $$\underline{\Aut}_{G,C}:A\mapsto \{f\in \Aut_A(G_A)\mid f|_{C_A}=\id_{C_A}\}$$ is affine of finite type and has maximal smooth closed $k$-subgroup $Z_{G,C}$ by \cite[2.4.1]{CGP15}. With this notation, we say $G$ is \emph{generalised standard} if there is a $4$-tuple $(G',k'/k,T',\Cs)$ such that $(G',k'/k)$ is a primitive pair, $T'$ a maximal torus of $G'$, $\Cs$  a commutative pseudo-reductive group, and there is a factorisation 
\begin{equation}\mathscr{C}\stackrel{\phi}{\to} \Cs\stackrel{\psi}{\to}Z_{\mathscr{G},\C}=\R_{k'/k}(Z_{G',C'})\label{fact}\end{equation}
with $\mathscr{G}=\mathscr{D}(R_{k'/k}(G'))$, $C'=Z_{G'}(T')$,  and $\C=\mathscr{G}\cap \R_{k'/k}(C')$---a Cartan $k$-subgroup of $\mathscr{G}$---such that there is a $k$-isomorphism 
\begin{equation}(\mathscr{G}\rtimes \Cs)/\C\cong G\label{gisaquotient}\end{equation}
where $\C$ is anti-diagonally embedded as a central $k$-subgroup of $\mathscr{G}\rtimes \Cs$. 
In our situation, since we are able to assume that $\ker(i_G)=1$, we see that we do not need to consider possible factors as in \cref{nonredcase}. For \cref{nonredcase}, the reader can refer to \cite[Thm.~9.8.6]{CGP15} (or \cite{BRSS24} for a more elementary construction).

%We can therefore deduce that the Weil restrictions $G_i:=\R_{k'/k}(G_i')$ are all perfect, hence $\G=\prod_{i=1}^r G_i$. (For the perfectness in case (i), see \cite[1.3.4]{CGP15}; for case(ii)(b) the group is the universal smooth tame central extension of a perfect group, so is itself perfect by \cite[5.1.3]{CPClass}---see 8.1.1 and 7.3.3 of \emph{op.~cit.} for the definition; in the other cases, the groups involved are absolutely pseudo-simple, hence perfect, \cite[2.2.2, 8.2.3, 8.3.6]{CPClass}.) 

Since the formation of $K_C(\lambda)$ commutes with separable extension (\Cref{mfodsep}), we may as well assume $k=k_s$. In that case the group $G'$ is pseudo-split with absolutely pseudo-simple fibres $G_i'$ over the factor fields $k_i'$ of $k'\cong \prod_{1\leq i\leq r}k_i'$ and there is some algebraic data that determines the possibilities for each $G_i'$ exactly. In the following lemma we use this data to give a description of the Cartan subgroup $C_i':=Z_{G_i'}(T'\cap G_i')$ of $G_i'$, thence a Cartan subgroup of $\D(\R_{k_i'/k}(G_i'))$. Since everything in sight distributes over direct products, we lose nothing by assuming $r=1$ for ease of notation---so $k'$ is a field. Let $\Delta=\Delta_>\cup \Delta_<$ be the root system of $G'$, with $\Delta_>$ the long and $\Delta_<$ the short roots.
\begin{lemma}
  Suppose $k=k_s$ and $(G',k'/k)$ is a primitive pair of type (i) or (ii). Let $T'$ be a split maximal torus of $G'$ and $K$ the minimal field of definition of $\RR_u((G')_\bark)$. Then \cref{tabledata} describes a Cartan subgroup of $\R_{k'/k}(G')$.
\end{lemma}

By way of notation, recall that if $G''$ is a split simple $K$-group and $K/k'$ is a finite field extension, then $\widehat{G'}:=\R_{K/k'}(G'')$ is a pseudo-split pseudo-reductive $k'$-group with root groups isomorphic to $\R_{K/k'}(\Ga)$. Then in special characteristics one may replace some of the root groups of $\widehat{G'}$ with a vector group $\underline{V}$, where $\underline V$ corresponds to a $k'$-subspace $V$ of $K$. Furthermore, if $V$ is a $k$-subspace of a finite field extension $K$ of $k$, then $(V)^*_{K/k}$ denotes the Zariski closure in $\R_{K/k}(\Gm)$ of the ratios of non-zero elements of $V$. 

\begin{table}\small \begin{tabular}{l|l|l|l}
Case & Data & Explanation & Cartan of $\D(\R_{k'/k}(G'))$ \\\hline\hline
(i) & root system & & $\R_{k'/k}(T')$\\\hline

(ii)(a) & \parbox[t]{0.22\textwidth}{root system; $K$} & 
\parbox[t]{0.25\textwidth}{$(G')_a\cong \R_{K/k}(\Ga)$, $a$ long;\\ $(G')_a\cong \Ga$,  $a$ short} & $\R_{K/k}(T')$\\\hline

(ii)(b) & \parbox[t]{0.22\textwidth}{rank; $K$; \\ $k'$-subspace $V$ of $K$\\ such that $k'\langle V\rangle = K$} & 
\parbox[t]{0.25\textwidth}{$(G')_b=\underline{V}$, $b$ short; \\ $(G')_b\cong \Ga$, $b$ long} & 
\parbox[t]{0.37\textwidth}{$\left(\prod_{a\in\Delta_>} a^\vee (\R_{k'/k}(\Gm))\right) \\\phantom{something}\times (\R_{K/k}(b_K^\vee))(V_{K/k}^*)$\\where $\Delta_<=\{b\}$}\\\hline

(ii)(c) & \parbox[t]{0.22\textwidth}{rank; $K$; \\ $k'$-subspace $V_>$ of $K$,\\defining subfield\\ $K_>=k'\langle V_>\rangle$} & 
\parbox[t]{0.25\textwidth}{$(G')_b=\underline{V_>}$, $b$ long;\\ $(G')_b\cong \R_{K/k'}(\Ga)$, $b$ short} & \parbox[t]{0.37\textwidth}{$\left(\prod_{a\in\Delta_<} a^\vee (\R_{K/k}(\Gm))\right)$ \\ \phantom{something} $\times (\R_{K/k}(b_K^\vee))((V_>)_{K/k}^*)$ \\where $\Delta_>=\{b\}$}\\\hline

(ii)(d) & \parbox[t]{0.22\textwidth}{$K$; \\ $k'$-subspace $V_>$ of $K$,\\defining subfield\\ $K_>=k'\langle V_>\rangle$;\\
$K_>$-subspace $V_<$ of $K$\\ with $K=k'\langle V_<\rangle$} & 
\parbox[t]{0.25\textwidth}{$(G')_b=\underline{V_>}$, $b$ long;\\
$(G')_a=\underline{V_<}$, $a$ short} & 
\parbox[t]{0.32\textwidth}{$(V_>)_{K_>/k}^*\times (V_<)^*_{K/k}$}\\\hline
\end{tabular}\caption{Data describing $G'$ in cases (i) and (ii)\label{tabledata}}\end{table}
\begin{proof}
The Cartan $k'$-subgroup $C'=Z_{G'}(T')$ is described by \cite[3.4.1,\ 8.2.5,\ 3.2.7]{CPClass} and agrees with that in \cref{tabledata} on taking $k=k'$. In general, $\R_{k'/k}(C')$ is a Cartan subgroup of $\R_{k'/k}(G')$. In cases (i) and (ii)(a), the latter group is perfect by \cite[1.3.4, 8.1.2]{CGP15} and the result is immediate; so we may assume we are in one of the remaining three types and $p=2$. If $\rk(G')=1$, then $G'$ is non-standard and isomorphic to $H_{V,K/k'}$ for some $k'$-subspace $V$ of $K$; \cite[8.1.1]{CPClass}. Then $\D(\R_{k'/k}(G'))$ has Cartan subgroup $V^*_{K/k}$ as required by \cite[8.1.3]{CPClass} (with $\ell'=k$). If $\rk(G')=n>2$, then the $A_{n-1}$ subgroup generated by the long (resp.~short) root groups of $G'$ in case (ii)(b) (resp.~(ii)(c)) is isomorphic to $\SL_{n}$ (resp.~$\R_{K/k'}(\SL_{n})$) and so its Weil restriction is perfect. Hence, the long root (resp.~short root) factor of a Cartan subgroup of $\R_{k'/k}(G')$ survives after passing to the derived subgroup. This reduces the assertion to rank $1$, which we have established already---one may appeal to \cite[C.2.32]{CGP15} if so desired. Similar arguments apply when $\rk(G')=2$ and are omitted.
\end{proof}

With \cref{tabledata} in hand, we can specify $K_C(\lambda)$ just in terms of the data therein and the field $K_{\Cs}(\lambda_\Cs)$, where $\lambda_\Cs$ is the restriction of $\lambda$ to a maximal split torus $\Ts$ of $\Cs$---in other words, we have a complete understanding modulo the commutative case. Let us arrange that the isomorphism of $T'$ with some $(\Gm)^n\cong (\Gm)^{n_1}\times\dots\times (\Gm)^{n_r}$ is lined up with the descriptions in the table, so that $G_i'\cap T'$ is the factor $(\Gm)^{n_i}$ and its centraliser in $G_i'$ is the given Cartan subgroup.

 Fix $i$ and let $\Delta=\Delta_>\sqcup\Delta_<$ be a base for the root system of $G_i'$. Let $e$, $e_>$ and $e_<$ denote the largest exponents of $p$ dividing $\langle\lambda,a\rangle$ for every $a\in \Delta$, $\Delta_>$ and $\Delta_<$ respectively. Then with reference to \cref{tabledata} we define

\begin{equation}K_i:=\begin{cases}
  k(k_i')^{p^{e}}\text{ if $G_i'$ is as in case (i)};\\
  k(K)^{p^{e}}\text{ if $G_i'$ is as in case (ii)(a)};\\
  k(K)^{p^{e_<}}\text{ if $G_i'$ is as in case (ii)(b);}\\
  \text{the compositum }k(K_>)^{p^{e_>}}(K)^{p^{e_<}}\text{ if $G_i'$ is as in cases (ii)(c) or (ii)(d)}.
\end{cases}\label{theki}\end{equation}

We come to the main theorem of this section.
\begin{theorem}\label{fieldfromclass}
  Let $\KK$ denote the compositum of all the $K_i$ together with $K_\Cs(\lambda_\Cs)$. Then 
\begin{equation}K_C(\lambda)= \mathcal{K}\label{theequation}\end{equation} 
\end{theorem}

Let $K/k$ be purely inseparable, $U$ a $k$-subspace of $K$ such that $U$ generates $K$ as a subfield, and put $\mathscr{U}:=U^*_{K/k}\subseteq \R_{K/k}(\Gm)=:\mathscr{K}$. The proof of the theorem above will require an understanding of the representations of groups $\mathscr{U}:=U^*_{K/k}\subseteq \R_{K/k}(\Gm)=:\mathscr{K}$. Such groups $\mathscr{U}$ are rather mysterious: for example, their dimensions seem to be impossible to predict easily---see \cite[9.1.8--9.1.10]{CGP15}. Nonetheless, the following lemma explains that their representation theory is easy. Being pseudo-reductive of rank $1$, the simple modules of $\mathscr{K}$ and $\mathscr{U}$ up to isomorphism are denoted $L_\mathscr{K}(\lambda)$ and $L_\mathscr{U}(\lambda)$ by \cref{thm:fromBS}; here, $\lambda$ indicates the weight of a maximal split torus. It follows that $\Res^{\mathscr{K}}_\mathscr{U}(L_\mathscr{K}(\lambda))$ is an isotypic direct sum of copies of $L_\mathscr{U}(\lambda)$. In fact:
\begin{lemma}\label{theustaractsirred}
   The restriction $\Res^{\mathscr{K}}_\mathscr{U}(L_\mathscr{K}(\lambda))$ is irreducible.
\end{lemma}
\begin{proof}
  If $\lambda=0$ then $L_\mathscr{K}(\lambda)\cong k$ is the trivial module and the result is clear. So suppose $\lambda\neq 0$. 
  
From \cite[Thm.~5.8]{BS22}, the action of $\mathscr{K}$ on $L_\mathscr{K}(\lambda)$ factors through $\mathscr{K}\to \mathscr{K}; x\mapsto x^a$ followed by the $p^s$-power map $\mathscr{K}\to \mathscr{K}^{p^s}\cong \R_{K^{p^s}/k}(\Gm)$ followed by an action of $\mathscr{K}^{p^s}$ on its natural module $L_{\mathscr{K}^{p^s}}(1)$---and the latter identifies with the field $kK^{p^s}$.

Let $\{u_1,\dots,u_d\}\in U\setminus\{0\}$ be a set of generators for $K$ as a $k$-algebra. Scaling as in \cite[3.1.4, Proof]{CPClass}, it does no harm to assume $u_1=1$, so that the ratios $u_i/u_j$---which are all $k$-points of $\mathscr{U}$---also contain a set of generators of $K$ as a $k$-algebra. The minimal $k$-subalgebra of $kK^{p^s}$ containing $1$ and stable under the group generated by the ratios $u_i^\lambda$ is the same as the subalgebra stable under $\mathscr{K}^{p^s}$; this is the whole of $K^{p^s}$ as required.
\end{proof}

\begin{proof}[Proof of \cref{fieldfromclass}]
Recall that we are working over $k=k_s$, and our assumption is that $G$ is of the form \cref{gisaquotient} and $V$ is a simple module for $G$. Set $D:=\End_G(V)$. We want to show that the minimal field of definition of $\Jac(D_\bark)$ is the compositum $\mathcal{K}$ of the fields referenced by the theorem. 
First observe that $V$ lifts to a simple module for the pseudo-split pseudo-reductive  semidirect product $\GG:=\G\rtimes\Cs$, through the quotient map $\G\rtimes\Cs\to G\cong(\G\rtimes\Cs)/\C$, by letting the central antidiagonal $\C$ in $\GG$ act trivially. We work with the Cartan subalgebra $\CCC:=\C\times\Cs$, which is the centraliser of the maximal split torus $\TT:=T'\times \Ts$ where $T'$ is the canonical maximal split torus in $\R_{k'/k}(T')$. (Of course the product is direct since $\CCC$ is commutative.) This surjects onto $C$ with kernel $\C$, where $\TT$ maps onto $T$. 

Now for any $\lambda\in X(T)$ we get a corresponding lift $\widehat{\lambda}\in X(\TT)$, and so we get an isomorphism $V\cong L_{\G}(\widehat\lambda)$. Evidently $\End_G(V)=\End_{\G\rtimes\Cs}(V)$ and so these algebras equally identify with both $K_C(\lambda)$ and $K_\CCC(\lambda)$. Hence we need only show $K_\CCC(\lambda)\cong\KK$. Since $K_\CCC(\lambda)$ is by definition identified with the high weight space of $L_\GG(\widehat{\lambda})$ it suffices to show that this is the compositum $\KK$ as claimed. 

From \cref{prop:splitend} (for example) one sees that $V|_{G_i}$ is isotypic and semisimple; indeed it is a direct sum of copies of $L_G(\widehat{\lambda_i})$, where $\lambda_i:=\widehat{\lambda}|_{T_i}$. The  endomorphism algebra over $G_i$ is therefore the field $L_{C_i}(\lambda_i)$ and we wish to see that this identifies with the field $K_i$ in \cref{theki}, which we now do case-by-case. 

In case (i) $C_i=\R_{k_i'/k}(T_i')$ and the statement that $L_{C_i}(\lambda_i)\cong K_i$ is \cite[Thm.~5.8]{BS22}. The same result also deals with case (ii)(a).  Then we have $p=2$. We treat case (ii)(d), the others being similar. 

% Note that if the ratios of non-zero elements of $U$ generate $K$ over $k$ as a subalgebra of $K$ (that is, in the notation of \cref{tabledata}, if $k\langle U\rangle = K$), then the $k$-representations of $(U)^*_{K/k}$ and $\R_{K/k}(\Gm)$ are identical. For any irreducible representation of the latter, say $L_{\R_{K/k}(\Gm)}(\mu)$, identifies with the subfield $K(\mu)=kK^{p^{v_p(\mu)}}$ of $K$; so since $1\in K(\mu)$ generates it as a $k$-vector space over $(K^{p^{v_p(\mu)}})^\times$, then so must $(U)^*_{K/k}$, whose action factors through the faithful action of $(U^{p^{v_p(\mu)}})^*_{K/k}$---the $k$-points of the latter generate $K^{p^{v_p(\mu)}}$ over $k$ because the $p$-power map commutes with multiplication and addition. 

% Applying these observations to the Cartan subgroup $(V_>)^*_{K_>/k}\times (V_<)^*_{K/k}$ of $G_i$, we find we may as well work with $\R_{K_>/k}(\Gm)\times \R_{K/k}(\Gm)$ whence we can appeal again to \cite[Thm.~5.8]{BS22}. 
By \cref{theustaractsirred} the Cartan subgroup $(V_>)^*_{K_>/k}\times (V_<)^*_{K/k}$ of $G_i$ has the same irreducible representations as $\R_{K_>/k}(\Gm)\times \R_{K/k}(\Gm)$, whence we can appeal again to \cite[Thm.~5.8]{BS22}. Now apply \cref{directproduct} (inductively) to the product $\GG=\prod G_i\rtimes\Cs$ to see that $K_C(\lambda)\cong\End_\GG(V)$ is the compositum $\KK$ of the fields $K_i$ together with $K_\Cs(\lambda_\Cs)$ as required.
\end{proof}

The Conrad--Prasad structure theorem also gives a refinement to our dimension formula in \cref{thm:fromBS}(iv). As in the proof of \cref{fieldfromclass}, we may lift the action of a pseudo-split $G$ on $L_G(\lambda)$ to that of $\GG=\G\rtimes \Cs$, where we have accordingly a decomposition $\TT=T'\times \Ts$ of a split maximal torus of $\GG$. Let $M\supseteq \TT$ denote a split Levi subgroup of $\GG$. Then $M\cap \Cs=\Ts$ and let $M_i:=M\cap G_i$ with $T_i:=T\cap G_i$ a corresponding maximal split torus; lastly, set $\lambda_i$ (resp.~$\lambda_\Ts$) the restriction of $\lambda$ to $T_i$ (resp.~$\Ts$). Since the $M_i$ are absolutely simple and simply connected, we have $M\cong M_1\times\dots\times M_r\times \Ts$, and $\End_{M_i}(L_{M_i}(\lambda_i))=k$ is a trivial $M$-module. Using \cref{directproduct} and \cref{redtofactors} it follows that $L_M(\lambda)\cong L_{M_1}(\lambda_1)\otimes\dots\otimes L_{M_r}(\lambda_r)\otimes L_{\Ts}(\lambda_{\Ts})$. Since  $L_{\Ts}(\lambda_{\Ts})$ is just a $1$-dimensional weight module $k_{\lambda_\Ts}$ with $t\cdot x=\lambda(t)x$ for $t\in \Ts(k)$, it follows that $\dim L_M(\lambda)=\prod_{i=1}^r \ell_i$ where $\ell_i=\dim L_{M_i}(\lambda_i)$. (All of this is  well-known.)
  \begin{cor}\label{cor:dimforsimples}
 We have \begin{equation}L_\GG(\lambda)|_M\cong \left(L_{M_1}(\lambda_1)\otimes\dots\otimes L_{M_r}(\lambda_r)\otimes L_{\Ts}(\lambda_\Ts)\right)^{\oplus \dim \mathcal{K}}.\label{tensorproddec}\end{equation}
Hence $\dim L_\GG(\lambda)=\prod \ell_i\cdot [\mathcal{K}:k]$, where $\mathcal{K}$ is the compositum in \cref{fieldfromclass}.

If $G$ is not necessarily pseudo-split, we have $\dim V=\dim_{k_s} L_{G_{k_s}}(\lambda)\cdot d\cdot |\Lambda|$, where $\Lambda$ is the orbit of $\Gal(k_s/k)$ on $\lambda$, some composition factor $L_{G_{k_s}}(\lambda)$ of $V_{k_s}$ occurs with multiplicity $d>0$, and where $\dim_{k_s}L_{G_{k_s}}(\lambda)$ can be deduced from the previous formula.

Furthermore, $\End_\GG(L_\GG(\lambda))\cong \mathcal{K}$, so that $\dim\End_\GG(L_\GG(\lambda))=[\mathcal{K}:k]$. In the non-pseudo split case, with notation of \cref{eq:ksdecomp}, we have $\dim \End_G(V)= d^2\cdot |\Lambda|\cdot [K_1:k_s]$.
\end{cor} 
\begin{proof}
For the first statement, note the description of $L_M(\lambda)$ has already been established. Applying the formula from \cref{thm:fromBS}(iv) gives $\dim L_\GG(\lambda)=\dim L_M(\lambda)\cdot \dim L_C(\lambda)$; but $L_C(\lambda)\cong K_C(\lambda)\cong \KK$ by \cref{fieldfromclass}.

The second statement is immediate.
\end{proof}

{\bf{Acknowledgement:}} Many thanks to Brian Conrad for spotting an error in the first version of this paper, as well as helpful clarifications about the main theorem of \cite{CPClass}. The second author is supported by the Leverhulme Trust Research Project Grant number RPG-2021-080. 

{\footnotesize
\bibliographystyle{amsalpha}
\bibliography{bib}}

\end{document}



\providecommand{\bysame}{\leavevmode\hbox to3em{\hrulefill}\thinspace}
\providecommand{\MR}{\relax\ifhmode\unskip\space\fi MR }
% \MRhref is called by the amsart/book/proc definition of \MR.
\providecommand{\MRhref}[2]{%
  \href{http://www.ams.org/mathscinet-getitem?mr=#1}{#2}
}
\providecommand{\href}[2]{#2}
\begin{thebibliography}{WITO69}

\bibitem[AK11]{And11}
Henning~Haahr Andersen and Masaharu Kaneda, \emph{Rigidity of tilting modules}, Mosc. Math. J. \textbf{11} (2011), no.~1, 1--39, 181. \MR{2808210}

\bibitem[ANT44]{ANT}
Emil Artin, Cecil~J. Nesbitt, and Robert~M. Thrall, \emph{Rings with {M}inimum {C}ondition}, University of Michigan Publications in Mathematics, vol. no. 1, University of Michigan Press, Ann Arbor, MI, 1944. \MR{10543}

\bibitem[BDM11]{BDM}
C.~Bowman, S.~R. Doty, and S.~Martin, \emph{Decomposition of tensor products of modular irreducible representations for {${\rm SL}_3$}}, Int. Electron. J. Algebra \textbf{9} (2011), 177--219, With an appendix by C. M. Ringel. \MR{2753767}

\bibitem[Ben98]{Ben98}
D.~J. Benson, \emph{Representations and cohomology. {I}}, second ed., Cambridge Studies in Advanced Mathematics, vol.~30, Cambridge University Press, Cambridge, 1998, Basic representation theory of finite groups and associative algebras. \MR{1644252}

\bibitem[BK02]{BK02}
S.~Bouchiba and S.~Kabbaj, \emph{Tensor products of {C}ohen-{M}acaulay rings: solution to a problem of {G}rothendieck}, J. Algebra \textbf{252} (2002), no.~1, 65--73. \MR{1922384}

\bibitem[Bor91]{Bor91}
Armand Borel, \emph{Linear algebraic groups}, second ed., Graduate Texts in Mathematics, vol. 126, Springer-Verlag, New York, 1991. \MR{1102012}

\bibitem[Bou22]{BourbAlg8}
N.~Bourbaki, \emph{Elements of {M}athematics. {A}lgebra. {C}hapter 8}, Springer, Cham, [2022] \copyright2022, English translation of [3027127], Translated by Reinie Ern\'e. \MR{4696567}

\bibitem[BR19]{BR19}
Dave Benson and Zinovy Reichstein, \emph{Fields of definition for representations of associative algebras}, Proc. Edinb. Math. Soc. (2) \textbf{62} (2019), no.~1, 291--304. \MR{3938833}

\bibitem[BRSS24]{BRSS24}
Michael Bate, Gerhard R{\"o}hrle, Damian Sercombe, and David~I. Stewart, \emph{A construction of pseudo-reductive groups with non-reduced root systems}, Transformation Groups (2024).

\bibitem[BS22]{BS22}
Michael Bate and David~I. Stewart, \emph{Irreducible modules for pseudo-reductive groups}, J. Eur. Math. Soc. (JEMS) \textbf{24} (2022), no.~7, 2533--2553. \MR{4413772}

\bibitem[CGP15]{CGP15}
Brian Conrad, Ofer Gabber, and Gopal Prasad, \emph{Pseudo-reductive groups}, second ed., New Mathematical Monographs, vol.~26, Cambridge University Press, Cambridge, 2015. \MR{3362817}

\bibitem[CP16]{CPClass}
Brian Conrad and Gopal Prasad, \emph{Classification of pseudo-reductive groups.}, Ann. Math. Stud., vol. 191, Princeton, NJ: Princeton University Press, 2016 (English). \MR{3379926}

\bibitem[CP17]{CPSurv}
\bysame, \emph{Structure and classification of pseudo-reductive groups}, Algebraic groups: structure and actions, Proc. Sympos. Pure Math., vol.~94, Amer. Math. Soc., Providence, RI, 2017, pp.~127--276. \MR{3645069}

\bibitem[Eis95]{Eisenbud}
David Eisenbud, \emph{Commutative algebra}, Graduate Texts in Mathematics, vol. 150, Springer-Verlag, New York, 1995, With a view toward algebraic geometry. \MR{1322960}

\bibitem[GP11]{SGA3}
Philippe Gille and Patrick Polo (eds.), \emph{Sch\'{e}mas en groupes ({SGA} 3). {T}ome {III}. {S}tructure des sch\'{e}mas en groupes r\'{e}ductifs}, Documents Math\'{e}matiques (Paris) [Mathematical Documents (Paris)], vol.~8, Soci\'{e}t\'{e} Math\'{e}matique de France, Paris, 2011, S\'{e}minaire de G\'{e}om\'{e}trie Alg\'{e}brique du Bois Marie 1962--64. [Algebraic Geometry Seminar of Bois Marie 1962--64], A seminar directed by M. Demazure and A. Grothendieck with the collaboration of M. Artin, J.-E. Bertin, P. Gabriel, M. Raynaud and J-P. Serre, Revised and annotated edition of the 1970 French original. \MR{2867622}

\bibitem[Gro65]{EGAIV65}
A.~Grothendieck, \emph{\'{E}l\'{e}ments de g\'{e}om\'{e}trie alg\'{e}brique. {IV}. \'{E}tude locale des sch\'{e}mas et des morphismes de sch\'{e}mas. {II}}, Inst. Hautes \'{E}tudes Sci. Publ. Math. (1965), no.~24, 231. \MR{199181}

\bibitem[Han04]{Han04}
Timo Hanke, \emph{An explicit example of a noncrossed product division algebra}, Math. Nachr. \textbf{271} (2004), 51--68. \MR{2068883}

\bibitem[Haz17]{Haz17}
Amit Hazi, \emph{Radically filtered quasi-hereditary algebras and rigidity of tilting modules}, Math. Proc. Cambridge Philos. Soc. \textbf{163} (2017), no.~2, 265--288. \MR{3682630}

\bibitem[Iar94]{Iar94}
Anthony~A. Iarrobino, \emph{Associated graded algebra of a {G}orenstein {A}rtin algebra}, Mem. Amer. Math. Soc. \textbf{107} (1994), no.~514, viii+115. \MR{1184062}

\bibitem[Jan03]{Jan03}
J.~C. Jantzen, \emph{Representations of algebraic groups}, second ed., Mathematical Surveys and Monographs, vol. 107, American Mathematical Society, Providence, RI, 2003. \MR{MR2015057 (2004h:20061)}

\bibitem[Lam99]{Lam2}
T.~Y. Lam, \emph{Lectures on modules and rings}, Graduate Texts in Mathematics, vol. 189, Springer-Verlag, New York, 1999. \MR{1653294}

\bibitem[Lam01]{Lam1}
\bysame, \emph{A first course in noncommutative rings}, second ed., Graduate Texts in Mathematics, vol. 131, Springer-Verlag, New York, 2001. \MR{1838439}

\bibitem[Mac94]{Mac}
F.~S. Macaulay, \emph{The algebraic theory of modular systems}, Cambridge Mathematical Library, Cambridge University Press, Cambridge, 1994, Revised reprint of the 1916 original, With an introduction by Paul Roberts. \MR{1281612}

\bibitem[Mil17]{Milne17}
J.~S. Milne, \emph{Algebraic groups}, Cambridge Studies in Advanced Mathematics, vol. 170, Cambridge University Press, Cambridge, 2017.

\bibitem[Oes84]{Oes84}
Joseph Oesterl{\'e}, \emph{Nombres de {T}amagawa et groupes unipotents en caract{\'e}ristique {$p$}}, Invent. Math. \textbf{78} (1984), no.~1, 13--88. \MR{762353}

\bibitem[Ras71]{rasala}
Richard Rasala, \emph{Inseparable splitting theory}, Trans. Amer. Math. Soc. \textbf{162} (1971), 411--448. \MR{284421}

\bibitem[Sha77]{Sha77}
Rodney~Y. Sharp, \emph{The dimension of the tensor product of two field extensions}, Bull. London Math. Soc. \textbf{9} (1977), no.~1, 42--48. \MR{437510}

\bibitem[Tit71]{Tit71}
J.~Tits, \emph{Repr{\'e}sentations lin{\'e}aires irr{\'e}ductibles d'un groupe r{\'e}ductif sur un corps quelconque}, J. Reine Angew. Math. \textbf{247} (1971), 196--220. \MR{0277536}

\bibitem[vdK21]{vdk21}
Wilberd van~der Kallen, \emph{Reductivity properties over an affine base}, Indag. Math. (N.S.) \textbf{32} (2021), no.~5, 961--967. \MR{4310008}

\bibitem[WITO69]{watanabe}
Kei-ichi Watanabe, Takeshi Ishikawa, Sadao Tachibana, and Kayo Otsuka, \emph{On tensor products of {G}orenstein rings}, J. Math. Kyoto Univ. \textbf{9} (1969), 413--423. \MR{257062}

\end{thebibliography}
